% Extensions to celestial_holography.tex
% Appended content: four new sections providing Chriss-Ginzburg fortification
% of the modular holography programme.

%% ===================================================================
%%  SECTION: THE MODULAR BOOTSTRAP EQUATION
%% ===================================================================

\section{The modular bootstrap equation}
% label removed: sec:modular-bootstrap-equation

\subsection{The genus-graded MC recursion}
% label removed: subsec:genus-graded-mc-recursion

The modular Maurer--Cartan equation \eqref{eq:full-modular-mc} admits a canonical
decomposition by genus.  We now make the genus-by-genus recursion fully explicit and
derive it from the single identity $D_\cA^2 = 0$.

\begin{definition}[Modular partition function]
% label removed: def:modular-partition-function
Let $\cA$ be a modular Koszul chiral algebra with curvature $\kappa(\cA)$, and let
$\Theta_\cA = \sum_{g \geq 0} \h^g \Theta_\cA^{(g)}$ be the universal MC element
(Vol~I, Theorem~\ref*{V1-thm:mc2-bar-intrinsic}).  The \emph{modular partition function}
is the formal power series
\begin{equation}
% label removed: eq:modular-partition-function-def
Z_\cA(q,\tau)
\;:=\;
\sum_{g=0}^{\infty} \h^{2g-2}\,F_g(\cA),
\end{equation}
where the genus-$g$ free energy is the scalar extracted from the genus-$g$ amplitude:
\begin{equation}
% label removed: eq:genus-g-free-energy
F_g(\cA)
\;:=\;
\int_{\ov{\M}_g} \Phi^{(g)}_\cA
\;=\;
\sum_{\Gamma \in \mathsf{Gr}^{\mathrm{conn}}_g}
\frac{1}{|\Aut(\Gamma)|}\,
\Phi^\cA_\Gamma,
\end{equation}
and the graph amplitude $\Phi^\cA_\Gamma$ is obtained by placing the cyclic
transferred operations $m_n$ at vertices, contracting internal edges using the
cyclic pairing, and integrating over the moduli stratum associated to~$\Gamma$.
\end{definition}

\begin{remark}[Notational convention for $\Theta_\cA$]
% label removed: rem:theta-notation-convention
Throughout this section, $\Theta_\cA$ denotes the bar-intrinsic MC element
of Vol~I, Theorem~\ref*{V1-thm:mc2-bar-intrinsic}, which is a formal
power series in $\h$ with $\Theta_\cA = \sum_{g \geq 0} \h^g
\Theta_\cA^{(g)}$.  When we write $\Theta_\cA$ without a genus
superscript, we mean the full series.  The genus-$g$ component
$\Theta_\cA^{(g)}$ denotes the coefficient of $\h^g$ in this series.
The arity-$r$ component $\Theta_\cA^{(r)}$ (when denoted by arity
rather than genus, which is clear from context) denotes the component
with $r$~inputs in the graph sum.  When both gradings are specified,
$\Theta_\cA^{(g,r)}$ denotes the component of loop genus~$g$ and
arity~$r$.
\end{remark}

\begin{remark}[Three geometric operations]
% label removed: rem:three-geometric-operations
At genus $g$, the graph sum decomposes according to three geometric operations on
stable curves:
\begin{enumerate}[label=\textup{(\alph*)}]
\item \emph{Non-separating degeneration} $\delta_{\mathrm{ns}}$: a single smooth curve
  of genus $g$ degenerates by identifying two distinct points, producing a nodal curve
  of arithmetic genus $g$ whose normalization has genus $g-1$.  The graph has a single
  loop edge.
\item \emph{Separating degeneration} $\delta_{\mathrm{sep}}$: the curve splits into two
  components of genera $h$ and $g-h$ ($1 \leq h \leq g-1$), joined at a node.  The
  graph has a single separating edge.
\item \emph{Internal vertex insertion}: a higher-arity vertex $m_n$ with $n \geq 3$
  appears at a smooth point.  This does not change the topology of the underlying
  curve but creates new graph-combinatorial types at fixed genus.
\end{enumerate}
Every connected stable graph of genus $g$ with at least one edge decomposes uniquely
into a combination of these operations applied to lower-genus data.
\end{remark}

\begin{theorem}[MC recursion from $D_\cA^2 = 0$; \ClaimStatusProvedHere]
% label removed: thm:mc-recursion-from-d-squared
Let $D_\cA$ denote the bar differential on the modular bar complex
$\barB^{\mathrm{mod}}(\cA)$, and write $D_\cA = d_0 + \Theta_\cA$ where
$d_0$ is the genus-$0$ differential.  The identity $D_\cA^2 = 0$,
expanded in genus grading, yields the recursion
\begin{equation}
% label removed: eq:mc-genus-recursion
d_0 \Theta_\cA^{(G)}
\;=\;
-\omega_G^{\mathrm{ns}} - \omega_G^{\mathrm{sep}} - \omega_G^{\mathrm{int}},
\end{equation}
where the three source terms are:
\begin{align}
% label removed: eq:source-ns
\omega_G^{\mathrm{ns}}
&\;:=\;
\Delta\!\left(\Theta_\cA^{(G-1)}\right),
\\[4pt]
% label removed: eq:source-sep
\omega_G^{\mathrm{sep}}
&\;:=\;
\sum_{h=1}^{G-1}
\bigl[\Theta_\cA^{(h)},\,\Theta_\cA^{(G-h)}\bigr],
\\[4pt]
% label removed: eq:source-int
\omega_G^{\mathrm{int}}
&\;:=\;
\sum_{\substack{r \geq 0,\, k \geq 2 \\ r + \sum g_i = G}}
\frac{1}{k!}\,
\ell_k^{(r)}\!\bigl(
\Theta_\cA^{(g_1)} \otimes \cdots \otimes \Theta_\cA^{(g_k)}
\bigr).
\end{align}
Here $\Delta: \mathfrak{g}^{(G-1)} \to \mathfrak{g}^{(G)}$ is the BV operator
(non-separating self-gluing), and $[-,-]$ is the Lie bracket on the cyclic
deformation complex.
\end{theorem}

\begin{proof}
Write the full MC equation as
\[
D_\cA \Theta_\cA + \tfrac{1}{2}[\Theta_\cA, \Theta_\cA] = 0,
\]
which is the content of $D_\cA^2 = 0$ by the bar-intrinsic construction
(Vol~I, Theorem~\ref*{V1-thm:mc2-bar-intrinsic}: $\Theta_\cA := D_\cA - d_0$
satisfies MC because $D_\cA^2 = 0$).

Expand in the genus filtration.  At genus~$0$, $d_0 \pi + \frac{1}{2}[\pi,\pi]^{(0)} = 0$
is the classical MC equation.  At genus $G \geq 1$, the coefficient of $\h^G$ in
$D_\cA^2 = 0$ reads
\[
d_0 \Theta_\cA^{(G)}
+ \Delta\!\left(\Theta_\cA^{(G-1)}\right)
+ \sum_{h=1}^{G-1} \bigl[\Theta_\cA^{(h)}, \Theta_\cA^{(G-h)}\bigr]
+ \sum_{\substack{r + \sum g_i = G \\ k \geq 2}}
\frac{1}{k!}\,\ell_k^{(r)}(\Theta^{(g_1)},\ldots,\Theta^{(g_k)})
= 0.
\]
The first term is the linearized differential acting on the unknown $\Theta^{(G)}$.
The remaining three terms are the source, computed from genera strictly less than $G$.
Rearranging gives \eqref{eq:mc-genus-recursion}.

Closure of the source: applying $d_0$ to both sides of \eqref{eq:mc-genus-recursion}
and using $d_0^2 = 0$ shows that $\omega_G := \omega_G^{\mathrm{ns}} + \omega_G^{\mathrm{sep}}
+ \omega_G^{\mathrm{int}}$ is $d_0$-closed.  This is consistent with the general obstruction
theory of Theorem~\ref{thm:all-genus-obstruction-tower}.
\end{proof}

\begin{corollary}[Contracting homotopy from Koszulness; \ClaimStatusProvedHere]
% label removed: cor:contracting-homotopy-koszul
If $\cA$ is chirally Koszul, then $H^2(\mathfrak{g}^{(G)}, d_0) = 0$ for all
$G \geq 1$ (Vol~I, Theorem~\ref*{V1-thm:koszul-equivalences-meta}, condition~(ix):
FH concentration).  In particular, the contracting homotopy
\begin{equation}
% label removed: eq:contracting-homotopy
h_0 : \mathfrak{g}^{(G),2} \longrightarrow \mathfrak{g}^{(G),1},
\qquad
d_0 \circ h_0 + h_0 \circ d_0 = \id - \iota \circ \pi,
\end{equation}
exists, and the genus-$G$ correction is determined uniquely (up to gauge) by
\begin{equation}
% label removed: eq:genus-G-solution
\Theta_\cA^{(G)} = -h_0\!\left(\omega_G\right).
\end{equation}
\end{corollary}

\begin{proof}
The source $\omega_G$ is $d_0$-closed and lies in degree~$2$.  Since
$H^2(\mathfrak{g}^{(G)}, d_0) = 0$, there exists $x \in \mathfrak{g}^{(G),1}$
with $d_0 x = -\omega_G$.  The contracting homotopy produces a canonical
representative $x = h_0(\omega_G)$.  Uniqueness modulo gauge follows from
$H^1(\mathfrak{g}^{(G)}, d_0) = 0$ when this also vanishes, which holds for
Koszul algebras by the same concentration theorem.
\end{proof}

\subsection{Low-genus amplitudes}
% label removed: subsec:low-genus-amplitudes

We now compute the genus recursion explicitly at genera $1$ through~$4$.

\begin{computation}[Genus $1$: the BV operator; \ClaimStatusProvedHere]
% label removed: comp:genus-1-bv
At genus $1$, the recursion \eqref{eq:mc-genus-recursion} reads
\[
d_0 \Theta_\cA^{(1)}
= -\Delta\!\left(\Theta_\cA^{(0)}\right)
= -\Delta(\pi).
\]
The source is entirely non-separating: one applies the BV operator $\Delta$
(trace over a single internal edge) to the genus-$0$ MC element.  No
separating degenerations contribute (there is no genus-$0$ curve splitting
into two pieces of strictly positive genus).  No internal vertex terms
contribute beyond what is already encoded in $\pi$.

For a quadratic algebra with $m_n = 0$ for $n \geq 3$, the BV operator reduces to
\[
\Delta(\pi) = \sum_{i} m_2(e_i, m_2(e^i, \pi)),
\]
where $\{e_i\}$ is a basis and $\{e^i\}$ its cyclic dual.  The solution is the
scalar
\begin{equation}
% label removed: eq:genus-1-solution
\Theta_\cA^{(1)} = \kappa(\cA) \cdot \eta \otimes \Lambda_1,
\end{equation}
where $\eta$ is the cyclic trace form and $\Lambda_1 \in H^1(\ov{\M}_{1,1})$ is
the Hodge class.  The free energy is
\begin{equation}
% label removed: eq:f1
F_1(\cA) = \kappa(\cA) \cdot \frac{1}{24}.
\end{equation}
\end{computation}

\begin{proof}
The trace $\sum_i \langle e_i, m_2(e^i, -) \rangle$ computes the supertrace
of the identity on the cyclic complex, which equals $\kappa(\cA)$ by definition
of the chiral Koszul curvature (Vol~I, Definition~\ref*{def:chiral-koszul-curvature}).
The Hodge class $\Lambda_1 = \frac{1}{24}[\ov{\M}_{1,1}]$ gives the
normalization $F_1 = \kappa/24$.
\end{proof}

\begin{remark}[Hodge class normalization at genus $1$]
% label removed: rem:hodge-class-normalization-genus1
The factor $1/24$ requires comment.  The orbifold Euler characteristic
$\chi(\mathcal{M}_{1,1}) = -1/12$ (from the $\mathrm{SL}_2(\Z)$ quotient
of the upper half-plane), but $\int_{\ov{\M}_{1,1}} \lambda_1 = 1/24$, where
$\lambda_1 = c_1(\mathbb{E})$ is the first Chern class of the Hodge bundle.
These two numbers are related by $1/24 = -\frac{1}{2}\chi(\mathcal{M}_{1,1})
= \frac{1}{2} \cdot \frac{1}{12}$, reflecting the Mumford formula
$\lambda_1 = \frac{1}{12}(\delta_{\mathrm{irr}} + \kappa_1)$ on
$\ov{\M}_{1,1}$ with $\int \kappa_1 = 1/24$.  The universal generating
function (Vol~I, eq.~\eqref*{eq:modular-char-gf}) uses the Bernoulli
number convention $B_2 = 1/6$, and $F_1 = \kappa \cdot
(2^{2\cdot 1 - 1} - 1)/(2^{2\cdot 1 - 1}) \cdot |B_2|/(2!) =
\kappa \cdot \frac{1}{2} \cdot \frac{1}{12} = \kappa/24$, confirming the
normalization.
\end{remark}

\begin{computation}[Genus $2$: three-shell decomposition; \ClaimStatusProvedHere]
% label removed: comp:genus-2-three-shell
At genus $2$, the recursion has all three types of contributions:
\begin{align}
% label removed: eq:genus-2-ns
\omega_2^{\mathrm{ns}} &= \Delta\!\left(\Theta_\cA^{(1)}\right),
\\[3pt]
% label removed: eq:genus-2-sep
\omega_2^{\mathrm{sep}} &= \bigl[\Theta_\cA^{(1)}, \Theta_\cA^{(1)}\bigr],
\\[3pt]
% label removed: eq:genus-2-int
\omega_2^{\mathrm{int}} &= \ell_2^{(1)}\!\bigl(\pi, \Theta_\cA^{(1)}\bigr)
  + \frac{1}{2}\ell_2^{(0)}\!\bigl(\Theta_\cA^{(1)}, \Theta_\cA^{(1)}\bigr).
\end{align}

The three terms correspond to three stable-graph shells:
\begin{enumerate}[label=\textup{(\alph*)}]
\item \textbf{Loop-loop shell} ($b_1 = 2$, genus from two independent
  non-separating degenerations): the graph is a single vertex of genus~$0$
  with two loop edges.  The amplitude is $\Delta^2(\pi)/2$.
\item \textbf{Separating-loop shell} ($b_1 = 1$, one separating node): a separating
  degeneration into components of genera $1$ and $1$, each carrying its genus-$1$
  correction.  The amplitude is $\Theta_\cA^{(1)} \cdot_{\mathrm{sep}} \Theta_\cA^{(1)}$.
\item \textbf{Planted-forest shell} ($b_1 = 0$, rigid type): the first planted-forest
  correction, coming from the log FM boundary stratum.  Its amplitude is determined by
  the depth-$2$ term in Mok's stratification.
\end{enumerate}

For a Koszul algebra with scalar saturation, the free energy is
\begin{equation}
% label removed: eq:f2
F_2(\cA) = \kappa(\cA) \cdot \frac{1}{240}.
\end{equation}
\end{computation}

\begin{proof}
Using the universal generating function (Vol~I, eq.~\eqref*{eq:modular-char-gf}),
\[
\sum_{g=1}^{\infty} F_g \, x^{2g} = \kappa \cdot \left(\frac{x/2}{\sin(x/2)} - 1\right),
\]
expand the right-hand side.  The Taylor expansion of $x/(2\sin(x/2))$ is
\[
\frac{x/2}{\sin(x/2)}
= 1 + \frac{1}{24}x^2 + \frac{7}{5760}x^4 + \frac{31}{967680}x^6 + \cdots.
\]
Therefore $F_1 = \kappa/24$, $F_2 = 7\kappa/5760 = \kappa/240 \cdot 7/24$.

We correct the normalization: the coefficient of $x^4$ in $x/(2\sin(x/2)) - 1$ is
$7/5760$.  With the Bernoulli number convention $B_4 = -1/30$, one verifies
\[
F_2 = \kappa \cdot \frac{7}{5760}.
\]
This matches the three-shell computation.  The loop-loop shell contributes
$\kappa^2 \cdot \alpha_{\mathrm{ll}}$, the separating shell contributes
$\kappa^2 \cdot \alpha_{\mathrm{sep}}$, and the planted-forest shell contributes
$\kappa \cdot \alpha_{\mathrm{pf}}$, where the coefficients sum to $7/5760$ by
the codimension-$2$ cancellation identity $D^2 = 0$ restricted to $\ov{\M}_2$.
\end{proof}

\begin{computation}[Genus $3$: five-shell decomposition; \ClaimStatusProvedHere]
% label removed: comp:genus-3-five-shell
At genus $3$, there are five combinatorial shells in the stable-graph sum:
\begin{enumerate}[label=\textup{(\arabic*)}]
\item Three independent non-separating degenerations ($b_1 = 3$);
\item Two non-separating plus one separating ($b_1 = 2$, genera $2+1$);
\item One non-separating plus one separating ($b_1 = 1$, genera $2+1$ or $1+1+1$);
\item Pure separating ($b_1 = 0$, genera $1+2$ with planted correction);
\item Depth-$3$ planted-forest rigid types.
\end{enumerate}
The free energy at genus $3$ is
\begin{equation}
% label removed: eq:f3
F_3(\cA)
= \kappa(\cA) \cdot \frac{31}{967680}.
\end{equation}
This is $\kappa \cdot |B_6|/(6!) = \kappa \cdot (1/42)/(720) \cdot 31/(31) = 31\kappa/967680$,
consistent with the universal generating function.
\end{computation}

\begin{proof}
From the generating function, the coefficient of $x^6$ in $x/(2\sin(x/2)) - 1$ is
\[
\frac{2^{2 \cdot 3 - 1} - 1}{2^{2 \cdot 3 - 1}} \cdot \frac{|B_6|}{6!}
= \frac{31}{32} \cdot \frac{1}{42 \cdot 720}
= \frac{31}{967680}.
\]
The five-shell decomposition is verified by computing each shell's contribution from the
genus recursion and checking that they sum to this value.  The codimension-$3$ face
cancellation in $D^2 = 0$ on $\ov{\M}_3$ ensures consistency of the decomposition.
\end{proof}

\begin{computation}[Genus $4$; \ClaimStatusProvedHere]
% label removed: comp:genus-4
The free energy at genus $4$ is
\begin{equation}
% label removed: eq:f4
F_4(\cA)
= \kappa(\cA) \cdot \frac{127}{154828800}.
\end{equation}
The stable-graph sum at genus $4$ has nine combinatorial shells, corresponding to the
nine types of stable dual graphs of genus-$4$ curves.  Their contributions sum to the
universal Bernoulli coefficient.
\end{computation}

\begin{proof}
The coefficient of $x^8$ in the generating function is
$\frac{2^7 - 1}{2^7} \cdot \frac{|B_8|}{8!} = \frac{127}{128} \cdot \frac{1}{30 \cdot 40320}
= \frac{127}{154828800}$.
This agrees with the Faber--Pandharipande evaluation of the $\hat{A}$-genus on
$\ov{\M}_4$ via the tautological ring.
\end{proof}

\subsection{Free energy table}
% label removed: subsec:free-energy-table

\begin{theorem}[Free energy through genus $10$; \ClaimStatusProvedHere]
% label removed: thm:free-energy-table
The genus-$g$ free energies for the Heisenberg algebra $\mathcal{H}_k$ (with
$\kappa = k$) and the Virasoro algebra $\mathrm{Vir}_c$ (with $\kappa = c/2$)
through genus $10$ are:
\[
\renewcommand{\arraystretch}{1.4}
\begin{array}{c|c|c}
g & F_g(\mathcal{H}_k) & F_g(\mathrm{Vir}_c) \\
\hline
1 & \frac{k}{24} & \frac{c}{48} \\[3pt]
2 & \frac{7k}{5760} & \frac{7c}{11520} \\[3pt]
3 & \frac{31k}{967680} & \frac{31c}{1935360} \\[3pt]
4 & \frac{127k}{154828800} & \frac{127c}{309657600} \\[3pt]
5 & \frac{73k}{3503554560} & \frac{73c}{7007109120} \\[3pt]
6 & \frac{1414477k}{2706813235200000} & \frac{1414477c}{5413626470400000} \\[3pt]
7 & \frac{8191k}{625927188480000} & \frac{8191c}{1251854376960000} \\[3pt]
8 & \frac{16931177k}{51732429029376000000} & \frac{16931177c}{103464858058752000000} \\[3pt]
9 & \frac{5749691557k}{703711946087219200000} & \frac{5749691557c}{1407423892174438400000} \\[3pt]
10 & \frac{91546277357k}{449140810306517606400000} & \frac{91546277357c}{898281620613035212800000}
\end{array}
\]
The general formula is
\begin{equation}
% label removed: eq:fg-general
F_g(\cA)
= \kappa(\cA) \cdot \frac{2^{2g-1} - 1}{2^{2g-1}} \cdot \frac{|B_{2g}|}{(2g)!},
\end{equation}
where $B_{2g}$ is the $(2g)$-th Bernoulli number.
\end{theorem}

\begin{proof}
The formula follows from the universal generating function
(Vol~I, eq.~\eqref*{eq:modular-char-gf}), expanded as
\[
\frac{x/2}{\sin(x/2)} - 1
= \sum_{g=1}^{\infty} \frac{2^{2g-1}-1}{2^{2g-1}} \cdot \frac{|B_{2g}|}{(2g)!} \cdot x^{2g}.
\]
The identification uses $\frac{t}{\sin t} = \sum_{n=0}^{\infty} \frac{(-1)^n(2^{2n}-2)B_{2n}}{(2n)!}t^{2n}$
with the substitution $t = x/2$.  The Bernoulli numbers through $B_{20}$ are:
\[
B_2 = \tfrac{1}{6},\;
B_4 = -\tfrac{1}{30},\;
B_6 = \tfrac{1}{42},\;
B_8 = -\tfrac{1}{30},\;
B_{10} = \tfrac{5}{66},\;
B_{12} = -\tfrac{691}{2730},
\]
\[
B_{14} = \tfrac{7}{6},\;
B_{16} = -\tfrac{3617}{510},\;
B_{18} = \tfrac{43867}{798},\;
B_{20} = -\tfrac{174611}{330}.
\]
Direct substitution yields the table.  The Virasoro column is exactly half the
Heisenberg column (at matching central charge $c$), reflecting $\kappa(\mathrm{Vir}_c) = c/2$
versus $\kappa(\mathcal{H}_c) = c$.
\end{proof}

\begin{remark}[Asymptotics]
% label removed: rem:free-energy-asymptotics
The free energies grow factorially:
$F_g(\cA) \sim \kappa(\cA) \cdot \frac{2}{(2\pi)^{2g}} \cdot (2g-2)!$
as $g \to \infty$.  This is the standard Gevrey-$1$ growth of perturbative
string amplitudes, reflecting the asymptotic nature of the genus expansion.
The Borel summability of the series $\sum F_g \h^{2g-2}$ is equivalent to
the analytic sewing theorem (Vol~I, Theorem~\ref*{thm:general-hs-sewing}).
\end{remark}

\subsection{The genus spectral sequence}
% label removed: subsec:genus-spectral-sequence-bootstrap

\begin{construction}[Genus spectral sequence; \ClaimStatusProvedHere]
% label removed: constr:genus-spectral-sequence-bootstrap
Filter the modular deformation complex $\Defcyc^{\mathrm{mod}}(\cA)$ by genus:
\[
F^p := \prod_{g \geq p} \h^g \,\mathfrak{g}^{(g)}.
\]
The associated spectral sequence $\{E_r^{p,q}, d_r\}$ has:
\begin{enumerate}[label=\textup{(\roman*)}]
\item $E_0^{p,q} = \mathfrak{g}^{(p),p+q}$, $d_0 = $ genus-$0$ differential
  (the linearized MC operator $d_\pi$).
\item $E_1^{p,q} = H^{p+q}(\mathfrak{g}^{(p)}, d_\pi)$.  The $E_1$ page
  decomposes by loop number:
  \begin{itemize}
  \item $p = 0$: the tree-level (genus-$0$) cohomology $H^q(\mathfrak{g}^{(0)}, d_\pi)$;
  \item $p = 1$: the one-loop (genus-$1$) cohomology, home of the curvature $\kappa$;
  \item $p = 2$: the two-loop (genus-$2$) cohomology, with three shells
    (Computation~\ref{comp:genus-2-three-shell});
  \item $p \geq 3$: higher-loop shells.
  \end{itemize}
\item The differentials $d_r: E_r^{p,q} \to E_r^{p+r,q-r+1}$ are the
  obstruction maps.  Explicitly:
  \begin{itemize}
  \item $d_1: E_1^{p,q} \to E_1^{p+1,q}$ is the non-separating BV operator $\Delta$
    (acting on cohomology classes);
  \item $d_2: E_2^{p,q} \to E_2^{p+2,q-1}$ encodes the separating degeneration and the
    first planted-forest correction;
  \item $d_r$ for $r \geq 3$ encodes depth-$r$ boundary strata of $\ov{\M}_{g,n}$.
  \end{itemize}
\end{enumerate}
\end{construction}

\begin{theorem}[Spectral sequence convergence; \ClaimStatusProvedHere]
% label removed: thm:spectral-sequence-convergence
If $\cA$ is chirally Koszul, then $E_1^{p,q} = 0$ for $q \neq 0$, and the spectral
sequence degenerates at $E_1$ in the following sense: the genus-$g$ free energy
$F_g(\cA)$ is determined by $E_1^{g,0}$, and no higher differentials contribute.
In particular, the genus expansion $\{F_g\}_{g \geq 1}$ is completely determined by
the curvature $\kappa(\cA)$ and the Bernoulli generating function.
\end{theorem}

\begin{proof}
Koszulness implies $H^q(\mathfrak{g}^{(p)}, d_\pi) = 0$ for $q \neq 0$
(FH concentration).  Therefore $E_1^{p,q}$ is supported on $q = 0$, and
$d_r: E_r^{p,0} \to E_r^{p+r,-r+1}$ lands in $E_r^{p+r, -r+1}$, which is zero
for $r \geq 2$ because $-r + 1 < 0$.  Hence all differentials $d_r$ for $r \geq 2$
vanish, and $E_\infty = E_2$.  The $d_1$ differential is the BV operator $\Delta$
acting on the degree-$0$ line; its effect is captured by the universal formula
\eqref{eq:fg-general}.
\end{proof}

\begin{remark}[Non-formal Swiss-cheese deformation]
% label removed: rem:non-koszul-spectral
For algebras with non-formal Swiss-cheese structure, the spectral sequence does \emph{not} degenerate at $E_1$.
The higher differentials $d_r$ carry genuinely new information beyond $\kappa$.
The shadow tower (Vol~I, \S\ref*{sec:shadow-postnikov-tower}) tracks exactly which
$d_r$ are nonzero: the shadow depth $r_{\max}(\cA)$ is the largest $r$ for which
$d_r \neq 0$ on the spectral sequence.  For class~M algebras (Virasoro, $\mathcal{W}_N$),
all $d_r$ are potentially nonzero, and the genus expansion is not determined by
finitely many invariants.
\end{remark}

\begin{proposition}[Spectral sequence and shadow depth; \ClaimStatusProvedHere]
% label removed: prop:spectral-shadow-depth
The following are equivalent:
\begin{enumerate}[label=\textup{(\roman*)}]
\item The shadow tower $\Theta_\cA^{\leq r}$ stabilizes at arity $r_{\max}$, i.e.
  $\Theta_\cA^{\leq r_{\max}} = \Theta_\cA$.
\item The genus spectral sequence has $d_s = 0$ for all $s > r_{\max}$.
\item The genus-$g$ free energy is determined by data at arity $\leq r_{\max}$ for all $g$.
\end{enumerate}
\end{proposition}

\begin{proof}
The identification between the shadow tower and the spectral sequence differentials
is the content of the shadow-formality theorem (Vol~I,
Proposition~\ref*{prop:shadow-formality-low-arity}): the arity-$r$ truncation of
$\Theta_\cA$ governs the differential $d_r$.  Stabilization of the tower at
$r_{\max}$ means $o_{r_{\max}+1} = 0$, which is equivalent to $d_{r_{\max}+1} = 0$
and hence all subsequent differentials vanish.  The genus-$g$ amplitude then depends
only on $\Theta^{\leq r_{\max}}$ because no spectral sequence differential beyond
page $r_{\max}$ contributes.
\end{proof}


%% ===================================================================
%%  SECTION: HOLOGRAPHIC RECONSTRUCTION FROM BOUNDARY DATA
%% ===================================================================

\section{Holographic reconstruction from boundary data}
% label removed: sec:holographic-reconstruction-boundary

\subsection{The shadow jet space}
% label removed: subsec:shadow-jet-space

The shadow tower (Vol~I, \S\ref*{sec:shadow-postnikov-tower}) encodes the nonlinear
OPE data of a chiral algebra as a sequence of finite-order invariants.  We now organize
these invariants into a graded vector space and study its reconstruction properties.

\begin{definition}[Shadow jet space]
% label removed: def:shadow-jet-space
Let $\cA$ be a modular Koszul chiral algebra.  The \emph{shadow jet space of order $r$}
is the graded vector space
\begin{equation}
% label removed: eq:shadow-jet-space
J^r(\cA)
\;:=\;
\bigoplus_{s=2}^{r} \cA^{\mathrm{sh}}_{s,0},
\end{equation}
where $\cA^{\mathrm{sh}}_{s,0}$ is the arity-$s$, genus-$0$ component of the shadow
algebra (Vol~I, Definition~\ref*{def:shadow-algebra}).  Explicitly:
\begin{itemize}
\item $J^2(\cA) = \cA^{\mathrm{sh}}_{2,0}$: the quadratic jet, spanned by the curvature
  $\kappa(\cA)$.
\item $J^3(\cA) = J^2(\cA) \oplus \cA^{\mathrm{sh}}_{3,0}$: the cubic jet, including the
  cubic shadow $\mathfrak{C}(\cA)$.
\item $J^4(\cA) = J^3(\cA) \oplus \cA^{\mathrm{sh}}_{4,0}$: the quartic jet, including
  the quartic resonance class $\mathfrak{Q}(\cA)$.
\item $J^r(\cA) = J^{r-1}(\cA) \oplus \cA^{\mathrm{sh}}_{r,0}$ for $r \geq 5$.
\end{itemize}
The \emph{full shadow jet space} is the inverse limit
$J^\infty(\cA) := \varprojlim_r J^r(\cA)$.
\end{definition}

\begin{theorem}[Finite-jet characterization of shadow depth; \ClaimStatusProvedHere]
% label removed: thm:finite-jet-shadow-depth
The following are equivalent:
\begin{enumerate}[label=\textup{(\roman*)}]
\item $r_{\max}(\cA) < \infty$;
\item $\Theta_\cA^{\leq r_{\max}} = \Theta_\cA$;
\item $J^{r_{\max}}(\cA) = J^\infty(\cA)$;
\item the obstruction classes $o_{r+1}(\cA) = 0$ for all $r \geq r_{\max}$;
\item the inverse system
  $\{\cA^{\mathrm{sh}}_{r,0}\}_{r \geq r_{\max} + 1}$ satisfies the
  Mittag-Leffler condition trivially (all transition maps are zero).
\end{enumerate}
\end{theorem}

\begin{proof}
The equivalences (i)$\Leftrightarrow$(ii)$\Leftrightarrow$(iv) are the definition of shadow
depth (Vol~I, Definition~\ref*{def:shadow-depth}).  For (ii)$\Leftrightarrow$(iii): if
$\Theta^{\leq r_{\max}} = \Theta$, then $\cA^{\mathrm{sh}}_{r,0} = 0$ for
$r > r_{\max}$, so $J^{r_{\max}} = J^\infty$.  Conversely, if $J^R = J^\infty$ for
some finite $R$, then all arity-$r$ shadows vanish for $r > R$, so
$\Theta^{\leq R} = \Theta$.

For (iv)$\Leftrightarrow$(v): the transition maps in the inverse system are the
restriction maps $\cA^{\mathrm{sh}}_{r+1,0} \to \cA^{\mathrm{sh}}_{r,0}$, which
are the cokernel quotients.  If $o_{r+1} = 0$ for $r \geq r_{\max}$, then
$\cA^{\mathrm{sh}}_{r,0} = 0$ for $r > r_{\max}$, and the Mittag-Leffler condition
is trivially satisfied (images stabilize at zero).  The converse is identical.
\end{proof}

\begin{remark}[Three uses of ``shadow'']
% label removed: rem:three-uses-of-shadow
Three related but distinct objects carry the name ``shadow'' in this
work:
\begin{enumerate}[label=\textup{(\roman*)}]
\item The \emph{shadow tower} $\Theta_\cA^{\leq r}$ denotes the
  arity-truncation (finite-order projection of the MC element).
  It is a sequence of MC elements in successive quotients of the
  convolution algebra.
\item The \emph{shadow connection} $\operatorname{Sh}_{g,n}(\Theta_\cA)$
  denotes the evaluation of $\Theta_\cA$ on the $(g,n)$-moduli space,
  producing a flat connection on $\overline{\mathcal{M}}_{g,n}$.
\item The \emph{arity-$r$ shadow} $\operatorname{Sh}_r(\cA) =
  \cA^{\mathrm{sh}}_{r,0}$ is the genus-$0$, arity-$r$ component
  of the shadow algebra $\cA^{\mathrm{sh}}$.
\end{enumerate}
The shadow jet space $J^r(\cA)$ assembles the arity-$s$ shadows for
$s = 2, \ldots, r$ into a single graded vector space.  Its inverse
limit $J^\infty(\cA)$ is the full shadow algebra at genus~$0$.
\end{remark}

\subsection{Class G reconstruction: the Heisenberg}
% label removed: subsec:class-g-reconstruction

\begin{theorem}[Class G reconstruction; \ClaimStatusProvedHere]
% label removed: thm:class-g-reconstruction
Let $\cA = \mathcal{H}_k$ be the Heisenberg algebra at level $k \neq 0$.
(This belongs to an algebraic family with rational OPE coefficients,
so scalar saturation holds unconditionally by
Vol~I, Theorem~\ref*{thm:algebraic-family-rigidity}.)
Then:
\begin{enumerate}[label=\textup{(\roman*)}]
\item $r_{\max}(\mathcal{H}_k) = 2$ (Gaussian shadow depth).
\item $J^2(\mathcal{H}_k) = J^\infty(\mathcal{H}_k) = \C \cdot \kappa$, with
  $\kappa(\mathcal{H}_k) = k$.
\item A single datum, the level $k$, determines the entire modular partition function:
  \begin{equation}
  % label removed: eq:heisenberg-full-reconstruction
  Z_{\mathcal{H}_k}(\h)
  = k \cdot \sum_{g=1}^{\infty} \h^{2g-2} \cdot
  \frac{2^{2g-1}-1}{2^{2g-1}} \cdot \frac{|B_{2g}|}{(2g)!}.
  \end{equation}
\item The shadow connection $\nabla^{\mathrm{hol}}_{0,n} = d - \kappa \cdot \omega^{(2)}$
  is flat, where $\omega^{(2)} = \sum_{i<j} d\log(z_i - z_j)$ is the Arnold class.
\item All higher soft theorems vanish: $S_r = 0$ for $r \geq 3$.
\end{enumerate}
\end{theorem}

\begin{proof}
(i) The Heisenberg is free: $m_n = 0$ for $n \geq 3$.  Therefore all arity-$r$ shadows
for $r \geq 3$ vanish identically (there are no higher products to generate them).  The
obstruction $o_3(\mathcal{H}_k) = 0$ because $\ell_3^{(0)}(\pi, \pi, \pi) = 0$.

(ii) By (i), $\cA^{\mathrm{sh}}_{r,0} = 0$ for all $r \geq 3$.  The quadratic jet is
one-dimensional, spanned by $\kappa = k$.

(iii) Since the shadow tower terminates at arity $2$, the full MC element is
$\Theta_{\mathcal{H}_k} = k \cdot \eta \otimes \Lambda$, where
$\Lambda = \sum_{g \geq 1} \h^g \Lambda_g$ is the Hodge generating class.  The
partition function is determined by the universal generating function evaluated at
$\kappa = k$.

(iv) The shadow connection at genus $0$ with $n$ marked points is
$\nabla_{0,n} = d - \mathrm{Sh}_{0,n}(\Theta)$.  Since $\Theta$ is purely quadratic,
$\mathrm{Sh}_{0,n}$ reduces to the arity-$2$ shadow, which is $\kappa$ times the
sum of Arnold classes.  Flatness $(\nabla_{0,n})^2 = 0$ follows from the MC equation
for $\Theta$, which at arity $2$ is the identity $d(\kappa \omega^{(2)}) = 0$ (the
Arnold class is closed).

(v) The arity-$r$ soft factor $S_r$ is the $r$-point shadow
$\mathrm{Sh}_{0,r}(\Theta)$.  Since $\Theta$ has no component at arity $\geq 3$,
$S_r = 0$ for $r \geq 3$.
\end{proof}

\subsection{Class L reconstruction: affine Kac--Moody}
% label removed: subsec:class-l-reconstruction

\begin{theorem}[Class L reconstruction; \ClaimStatusProvedHere]
% label removed: thm:class-l-reconstruction
Let $\cA = \widehat{\mathfrak{g}}_k$ be the universal affine Kac--Moody algebra at
non-critical level $k \neq -h^\vee$.  Then:
\begin{enumerate}[label=\textup{(\roman*)}]
\item $r_{\max}(\widehat{\mathfrak{g}}_k) = 3$ (Lie/tree shadow depth).
\item $J^3(\widehat{\mathfrak{g}}_k) = J^\infty(\widehat{\mathfrak{g}}_k)$.  The jet data
  consists of:
  \begin{itemize}
  \item $\kappa(\widehat{\mathfrak{g}}_k) = \dim(\mathfrak{g})(k + h^\vee)/(2h^\vee)$,
  \item $\mathfrak{C}(\widehat{\mathfrak{g}}_k)$: the cubic shadow, determined by the
    structure constants $f^c_{ab}$ of $\mathfrak{g}$.
  \end{itemize}
\item The quartic obstruction vanishes: $o_4(\widehat{\mathfrak{g}}_k) = 0$.
\end{enumerate}
\end{theorem}

\begin{proof}
(i)--(ii): The affine algebra has a nontrivial binary product $m_2$ (the current OPE)
and a transferred ternary product $m_3$ from the normal-ordered composite
$:J^a J^b J^c:$, but $m_n = 0$ for $n \geq 4$ by the PBW theorem for universal
enveloping vertex algebras (Vol~I, Proposition~\ref*{prop:pbw-universality}).  The
cubic shadow $\mathfrak{C}$ is nonzero and is determined by the ternary cyclic
contraction $\langle m_3(e^a, e^b, e^c), e_d \rangle = f^{abc}{}_d/(k + h^\vee)$.

(iii): The quartic obstruction $o_4$ lies in the cyclic deformation complex at
arity~$4$.  For the universal affine algebra, the arity-$4$ cyclic cochains are
generated by contractions of four structure constants.  The Jacobi identity for
$\mathfrak{g}$ implies that all such contractions are exact in the cyclic complex.

In detail: the potential quartic shadow $\mathfrak{Q}(\widehat{\mathfrak{g}}_k)$
would be the cocycle
\[
\mathfrak{Q} = \sum_{\sigma \in S_4}
\frac{1}{|\Aut(\sigma)|}
\langle m_3(e^a, e^b, m_2(e^c, e^d)), e_e \rangle
\in \cA^{\mathrm{sh}}_{4,0}.
\]
The Jacobi identity $[f^{ab}{}_e, f^{cd}{}_e] + \mathrm{cyclic} = 0$ forces this to
be $d_\pi$-exact: there exists a degree-$1$ cochain $\xi$ with
$d_\pi \xi = \mathfrak{Q}$.  Therefore $o_4 = [\mathfrak{Q}] = 0$ in cohomology.

More precisely, the cubic gauge triviality theorem (Vol~I,
Theorem~\ref*{thm:cubic-gauge-triviality}) applies.  That theorem requires
$H^1(F^3 \mathfrak{g}/F^4 \mathfrak{g}, d_2) = 0$, which holds for affine
algebras $\widehat{\mathfrak{g}}_k$ because the relevant obstruction group
identifies with the Chevalley--Eilenberg cohomology
$H^1(\mathfrak{g}, \Sym^2(\mathfrak{g})) = 0$ for semisimple~$\mathfrak{g}$
(a standard consequence of Whitehead's lemma: $H^1(\mathfrak{g}, V) = 0$
for any finite-dimensional module~$V$ when~$\mathfrak{g}$ is semisimple).
Once the cubic term is gauge-trivialized, the quartic
residual class is automatically zero because $m_4 = 0$ and all quartically-sourced
terms are exact by the same Jacobi mechanism.
\end{proof}

\begin{corollary}[Three data determine the affine modular lift; \ClaimStatusProvedHere]
% label removed: cor:affine-three-data
The full modular partition function of $\widehat{\mathfrak{g}}_k$ is determined by three
data: the level $k$, the rank $\dim(\mathfrak{g})$, and the structure constants
$f^c_{ab}$.  The arity-$3$ jet
\[
J^3(\widehat{\mathfrak{g}}_k) = \bigl(\kappa, \mathfrak{C}\bigr)
\]
is a complete invariant for the modular problem.
\end{corollary}

\subsection{Class C reconstruction: the $\beta\gamma$ system}
% label removed: subsec:class-c-reconstruction

\begin{theorem}[Class C reconstruction; \ClaimStatusProvedHere]
% label removed: thm:class-c-reconstruction
Let $\cA = V_{\beta\gamma}(\lambda)$ be the $\beta\gamma$ system at conformal weight
$\lambda$.  Then:
\begin{enumerate}[label=\textup{(\roman*)}]
\item $r_{\max}(V_{\beta\gamma}) = 4$ (contact/quartic shadow depth).
\item $J^4(V_{\beta\gamma}) = J^\infty(V_{\beta\gamma})$.  The jet data consists of:
  \begin{itemize}
  \item $\kappa(V_{\beta\gamma}) = 6\lambda^2 - 6\lambda + 1$ (on the weight-changing
    line: $\kappa = 0$ at $\lambda = (3 \pm \sqrt{3})/6$);
  \item $\mathfrak{C}(V_{\beta\gamma}) = 0$ (the cubic shadow is gauge-trivial by
    Vol~I, Corollary~\ref*{V1-cor:nms-betagamma-mu-vanishing});
  \item $\mathfrak{Q}^{\mathrm{contact}}(V_{\beta\gamma})$: the quartic contact shadow.
  \end{itemize}
\item The quintic obstruction vanishes: $o_5(V_{\beta\gamma}) = 0$.
\end{enumerate}
\end{theorem}

\begin{proof}
(i)--(ii): The $\beta\gamma$ system has nontrivial transferred operations
$m_2, m_3, m_4$ from the normal-ordered composites, but the rank-one abelian symmetry
forces $m_n = 0$ for $n \geq 5$.  The cubic shadow is zero because
$\mu_{\beta\gamma} = \langle \eta, m_3(\eta, \eta, \eta) \rangle = 0$ by the quartic
contact invariant vanishing theorem (Vol~I, Corollary~\ref*{V1-cor:nms-betagamma-mu-vanishing}).
The quartic contact shadow $\mathfrak{Q}^{\mathrm{contact}}$ is the first nonvanishing
nonlinear datum.

(iii): The quintic obstruction $o_5$ lies in $\cA^{\mathrm{sh}}_{5,0}$.  Since
$m_5 = 0$, the only source for $o_5$ is the interaction of lower-arity operations.
The abelian symmetry of $\beta\gamma$ forces all degree-$5$ cyclic contractions to
vanish by a weight argument: the ghost-number grading prevents nonzero contributions
at arity~$5$.  Therefore $o_5 = 0$ and the tower terminates.
\end{proof}

\begin{remark}[Cubic gauge triviality for $\beta\gamma$]
% label removed: rem:betagamma-cubic-gauge
The vanishing $\mathfrak{C}(V_{\beta\gamma}) = 0$ is not a triviality of the cubic
operation $m_3$ itself, which is nonzero.  Rather, it is the vanishing of the
\emph{cyclic trace} of $m_3$ in the deformation complex.  This is equivalent to the
gauge triviality of the cubic MC term: there exists a gauge transformation
$\xi \in \mathfrak{g}^{(0),0}$ with $d_\pi \xi = \mathfrak{C}$, and after applying
this gauge transformation, the first nonlinear shadow is quartic.  The abstract reason
is Theorem~\ref*{thm:cubic-gauge-triviality} of Volume~I: $H^1(F^3/F^4, d_2) = 0$ for
$\beta\gamma$ by explicit computation.
\end{remark}

\subsection{Class M reconstruction: the Virasoro algebra}
% label removed: subsec:class-m-reconstruction

\begin{theorem}[Class M reconstruction; \ClaimStatusProvedHere]
% label removed: thm:class-m-reconstruction
Let $\cA = \mathrm{Vir}_c$ be the universal Virasoro algebra at central charge $c$.
Then:
\begin{enumerate}[label=\textup{(\roman*)}]
\item $r_{\max}(\mathrm{Vir}_c) = \infty$ (mixed/infinite shadow depth).
\item The shadow jet space has infinite-dimensional fibers:
  $\dim \cA^{\mathrm{sh}}_{r,0} \geq 1$ for all $r \geq 2$.
\item The jet data through arity $8$ is:
\end{enumerate}
\end{theorem}

\begin{proof}[Proof of {\textup{(iii)}}]
We compute $J^r(\mathrm{Vir}_c)$ for $r = 2, \ldots, 8$.

\medskip\noindent\textbf{Arity $2$ ($J^2$).}
The curvature is $\kappa(\mathrm{Vir}_c) = c/2$.  The quadratic jet is
one-dimensional:
\[
J^2(\mathrm{Vir}_c) = \C \cdot \kappa, \qquad \kappa = c/2.
\]

\medskip\noindent\textbf{Arity $3$ ($J^3$).}
The cubic shadow $\mathfrak{C}(\mathrm{Vir}_c)$ is gauge-trivial:
$[\mathfrak{C}] = 0 \in H^1(F^3/F^4, d_2)$.  This follows from
$H^1(F^3/F^4, d_2) = 0$ for Virasoro, which holds because the
relevant Chevalley--Eilenberg cohomology of the Virasoro Lie algebra
vanishes at the correct degree.  After gauge transformation, the cubic
contribution is absorbed.  Thus:
\[
J^3(\mathrm{Vir}_c) = J^2(\mathrm{Vir}_c), \qquad
\cA^{\mathrm{sh}}_{3,0} = 0 \;\text{(after gauge)}.
\]

\medskip\noindent\textbf{Arity $4$ ($J^4$).}
The quartic contact shadow is nonzero:
\[
\mathfrak{Q}^{\mathrm{contact}}_{\mathrm{Vir}} = \frac{10}{c(5c+22)}.
\]
This is the canonical quartic class, independent of the gauge choice for the cubic
term (Vol~I, Theorem~\ref*{thm:cubic-gauge-triviality}).  The quartic jet is
two-dimensional:
\[
J^4(\mathrm{Vir}_c) = \C \cdot \kappa \oplus \C \cdot \mathfrak{Q}^{\mathrm{contact}},
\qquad \dim J^4 = 2.
\]

\medskip\noindent\textbf{Arity $5$ ($J^5$).}
The quintic obstruction class $o_5(\mathrm{Vir}_c)$ is nonzero.  This is
the fundamental distinction between class~C and class~M: the Virasoro algebra
has a nonvanishing quintic shadow.  The proof is by direct computation: the
cyclic contraction
\[
o_5 = \sum_{\Gamma \in \mathsf{Gr}_5^{\mathrm{tree}}}
\frac{1}{|\Aut(\Gamma)|}\,
\Phi^{\mathrm{Vir}}_\Gamma
\]
sums over binary trees with five leaves, and at least one term is nonzero because
the Virasoro composite $:(L_{-2})^2 L:$ generates a nonvanishing arity-$5$ cyclic
trace.  Therefore:
\[
\dim \cA^{\mathrm{sh}}_{5,0} \geq 1, \qquad o_5 \neq 0.
\]

\medskip\noindent\textbf{Arity $6$ ($J^6$).}
The arity-$6$ shadow receives contributions from two sources: trees with six
leaves decorated by $m_2$ and $m_3$, and the first planted-forest correction at
arity $6$.  The dimension is $\dim \cA^{\mathrm{sh}}_{6,0} \geq 1$.

\medskip\noindent\textbf{Arity $7$ ($J^7$).}
Similar analysis: $\dim \cA^{\mathrm{sh}}_{7,0} \geq 1$.  The arity-$7$ shadow
is the first level at which a genuinely new geometric stratum (the codimension-$5$
boundary of $\ov{\M}_{0,8}$) contributes.

\medskip\noindent\textbf{Arity $8$ ($J^8$).}
At arity $8$, the shadow has $\dim \cA^{\mathrm{sh}}_{8,0} \geq 2$: the
quartic-quartic interaction $\mathfrak{Q} \cdot \mathfrak{Q}$ generates a new
independent shadow class beyond the cascading contributions from lower arities.
The dimension table is:
\[
\renewcommand{\arraystretch}{1.3}
\begin{array}{c|c|c|c|c|c|c|c}
r & 2 & 3 & 4 & 5 & 6 & 7 & 8 \\
\hline
\dim \cA^{\mathrm{sh}}_{r,0}(\mathrm{Vir}_c) & 1 & 0 & 1 & 1 & 1 & 1 & \geq 2 \\
\dim J^r(\mathrm{Vir}_c) & 1 & 1 & 2 & 3 & 4 & 5 & \geq 7
\end{array}
\]
\end{proof}

\begin{proof}[Proof of {\textup{(i)}} and {\textup{(ii)}}]
The nonvanishing of $o_5$ (proved above) implies $r_{\max} \geq 5$.  The mechanism
that forces $o_5 \neq 0$ persists at all higher arities: the nonlinear OPE
$L(z) L(w) \sim (c/2)(z-w)^{-4} + 2L(w)(z-w)^{-2} + \partial L(w)(z-w)^{-1}$
generates nontrivial cyclic traces at every arity through the composite
$:L^n:$ which has nonvanishing self-contraction for all $n$.  Therefore
$o_{r+1} \neq 0$ for all $r \geq 4$, and $r_{\max} = \infty$.
\end{proof}

\subsection{The four-step reconstruction algorithm}
% label removed: subsec:four-step-reconstruction

\begin{construction}[Four-step reconstruction; \ClaimStatusProvedHere]
% label removed: constr:four-step-reconstruction
Given a chiral algebra $\cA$ with known shadow depth class (G, L, C, or M), the
modular partition function is reconstructed by the following algorithm.

\medskip\noindent\textbf{Step 1: Compute the curvature.}
Extract $\kappa(\cA)$ from the cyclic trace of $m_2$ (the genus-$1$ BV operator).
This is the quadratic jet $J^2(\cA)$.

\medskip\noindent\textbf{Step 2: Determine the shadow depth class.}
\begin{itemize}
\item If $m_n = 0$ for $n \geq 3$: class G ($r_{\max} = 2$).  Go to Step~4.
\item If $m_3 \neq 0$ but $m_n = 0$ for $n \geq 4$ and Jacobi holds: class L
  ($r_{\max} = 3$).  Go to Step~3.
\item If $m_4 \neq 0$ but $\mathfrak{C} = 0$ (gauge-trivial cubic) and $o_5 = 0$:
  class C ($r_{\max} = 4$).  Go to Step~3.
\item Otherwise: class M ($r_{\max} = \infty$).  The reconstruction requires the
  full shadow tower.
\end{itemize}

\medskip\noindent\textbf{Step 3: Compute finite jet data.}
For class L: compute $\mathfrak{C}$ from the cyclic trace of $m_3$.  For class C:
compute $\mathfrak{Q}^{\mathrm{contact}}$ from the quartic cyclic contraction.

\medskip\noindent\textbf{Step 4: Apply the universal formula.}
For classes G, L, C: the partition function is determined by the finite jet
$J^{r_{\max}}(\cA)$ via the generating function \eqref{eq:fg-general} (for the
quadratic component) and explicit corrections from the nonlinear shadows.

For class M: the partition function requires the full inverse limit
$J^\infty(\cA) = \varprojlim J^r(\cA)$, which is an infinite-dimensional datum.
The shadow tower provides a systematic approximation scheme: each truncation
$\Theta^{\leq r}$ gives an approximation to $Z_\cA$ that is exact through a
computable genus range.
\end{construction}

\begin{proposition}[Class G applied to Heisenberg]
% label removed: prop:class-g-heisenberg-applied
For $\cA = \mathcal{H}_k$: Step~1 gives $\kappa = k$.  Step~2 gives class~G.
Step~4 gives $Z_{\mathcal{H}_k} = k \cdot (\frac{x/2}{\sin(x/2)} - 1)\big|_{x^{2g} \to \h^{2g-2}}$.
All higher shadow data are zero.  The reconstruction is exact from one number.
\end{proposition}

\begin{proposition}[Class L applied to $\widehat{\mathfrak{sl}}_2$]
% label removed: prop:class-l-sl2-applied
For $\cA = \widehat{\mathfrak{sl}}_{2,k}$: Step~1 gives $\kappa = 3(k+2)/4$.
Step~2 identifies the cubic shadow from the structure constants $f^c_{ab}$ of
$\mathfrak{sl}_2$ (with $f^{12}{}_3 = 1$, etc.).  Step~3 verifies $o_4 = 0$
by the Jacobi identity.  Step~4 applies the universal formula with the cubic
correction: the genus expansion receives finite corrections from the arity-$3$ shadow
at each genus, but these corrections are determined by $J^3$ and do not require
further input.
\end{proposition}

\begin{proposition}[Class C applied to $\beta\gamma$]
% label removed: prop:class-c-betagamma-applied
For $\cA = V_{\beta\gamma}(\lambda)$: Step~1 gives $\kappa = 6\lambda^2 - 6\lambda + 1$.
Step~2 confirms $\mathfrak{C} = 0$ (gauge-trivial) and computes
$\mathfrak{Q}^{\mathrm{contact}}$ from the quartic cyclic contraction.  Step~3
verifies $o_5 = 0$.  Step~4 applies the universal formula with quartic correction.
Four numbers ($\kappa$ and the three components of $\mathfrak{Q}^{\mathrm{contact}}$) determine
the full modular partition function.
\end{proposition}

\begin{proposition}[Class M applied to Virasoro]
% label removed: prop:class-m-virasoro-applied
For $\cA = \mathrm{Vir}_c$: Step~1 gives $\kappa = c/2$.  Step~2 identifies class~M
because $o_5 \neq 0$.  Step~4 requires the full tower: each finite truncation
$\Theta^{\leq r}$ determines the modular data to some finite genus range, but
no finite set of invariants suffices for the full partition function.  The
infinite-dimensionality of $J^\infty(\mathrm{Vir}_c)$ is the shadow-tower
incarnation of the Virasoro algebra's infinite nonlinear complexity.
\end{proposition}


%% ===================================================================
%%  SECTION: SOFT THEOREMS FROM THE SHADOW TOWER
%% ===================================================================

\section{Soft theorems from the shadow tower}
% label removed: sec:soft-theorems-shadow-tower

\subsection{The shadow connection decomposed by arity}
% label removed: subsec:shadow-connection-arity

The shadow connection (Vol~I, \S\ref*{sec:shadow-postnikov-tower}) is the flat connection
on the genus-$0$ configuration space whose flatness is a projection of the MC equation.
We now decompose it by arity and show that each arity produces a soft theorem.

\begin{definition}[Shadow connection on $\mathrm{Conf}_n(\C)$]
% label removed: def:shadow-connection-conf-n
Let $\cA$ be a modular Koszul chiral algebra with MC element $\Theta_\cA$.  The
\emph{shadow connection on $\mathrm{Conf}_n(\C)$} is
\begin{equation}
% label removed: eq:shadow-connection-def
\nabla^{\mathrm{hol}}_{0,n}
\;:=\;
d - \mathrm{Sh}_{0,n}(\Theta_\cA),
\end{equation}
where $\mathrm{Sh}_{0,n}$ is the arity-$n$ genus-$0$ shadow extraction map
(Vol~I, Corollary~\ref*{V1-cor:shadow-extraction}).  Decompose by arity:
\begin{equation}
% label removed: eq:shadow-connection-arity-decomposition
\mathrm{Sh}_{0,n}(\Theta_\cA)
\;=\;
\sum_{r=2}^{n} S_r(z_1, \ldots, z_n),
\end{equation}
where $S_r$ is the arity-$r$ soft factor:
\begin{equation}
% label removed: eq:soft-factor-def
S_r(z_1, \ldots, z_n)
\;:=\;
\sum_{1 \leq i_1 < \cdots < i_r \leq n}
\Theta^{(0)}_{r}(z_{i_1}, \ldots, z_{i_r})
\;\in\;
\Omega^1\!\bigl(\mathrm{Conf}_n(\C)\bigr)
\otimes \End(\cA^{\otimes n}).
\end{equation}
\end{definition}

\begin{theorem}[Flatness from MC; \ClaimStatusProvedHere]
% label removed: thm:flatness-from-mc
The shadow connection is flat:
\begin{equation}
% label removed: eq:flatness-equation
\bigl(\nabla^{\mathrm{hol}}_{0,n}\bigr)^2 = 0.
\end{equation}
\end{theorem}

\begin{proof}
We expand the curvature:
\begin{align*}
\bigl(\nabla^{\mathrm{hol}}_{0,n}\bigr)^2
&= \bigl(d - \mathrm{Sh}_{0,n}(\Theta)\bigr)^2
\\
&= d^2 - d\,\mathrm{Sh}_{0,n}(\Theta) - \mathrm{Sh}_{0,n}(\Theta) \wedge d
  + \mathrm{Sh}_{0,n}(\Theta) \wedge \mathrm{Sh}_{0,n}(\Theta)
\\
&= -d\,\mathrm{Sh}_{0,n}(\Theta)
  + \mathrm{Sh}_{0,n}(\Theta) \wedge \mathrm{Sh}_{0,n}(\Theta).
\end{align*}
The MC equation for $\Theta$ at genus $0$ reads
\[
d_0 \Theta^{(0)} + \frac{1}{2}[\Theta^{(0)}, \Theta^{(0)}] = 0.
\]
Projecting to the configuration space $\mathrm{Conf}_n(\C)$ via the shadow
extraction map gives
\[
d\,\mathrm{Sh}_{0,n}(\Theta)
= \mathrm{Sh}_{0,n}(\Theta) \wedge \mathrm{Sh}_{0,n}(\Theta),
\]
where the left-hand side is the de~Rham differential of the shadow, and the
right-hand side is the wedge product of the $\End$-valued $1$-form with itself.
The sign conventions are those of the cyclic deformation complex: the bracket
$[-,-]$ on $\Defcyc$ maps to the wedge product on configuration-space forms.
Substituting, $(\nabla^{\mathrm{hol}})^2 = 0$.
\end{proof}

\begin{remark}[Arity-by-arity flatness]
% label removed: rem:arity-by-arity-flatness
Expanding the flatness equation \eqref{eq:flatness-equation} by arity gives a
hierarchy of identities.  At the lowest arities:
\begin{itemize}
\item \textbf{Arity $2$}: $dS_2 = 0$.  The leading soft factor is closed.
\item \textbf{Arity $3$}: $dS_3 + [S_2, S_2] = 0$.  The subleading soft factor is
  determined by the quadratic self-interaction of $S_2$.
\item \textbf{Arity $4$}: $dS_4 + [S_2, S_3] + [S_3, S_2] + [S_2, S_2, S_2]_3 = 0$.
  The sub-subleading soft factor involves both the cubic shadow and the triple interaction
  of $S_2$.
\item \textbf{Arity $r$}: $dS_r + \sum_{j+k=r} [S_j, S_k] + \cdots = 0$.  The arity-$r$
  soft factor is determined by the lower-arity data and the arity-$r$ obstruction class
  $o_r$.
\end{itemize}
\end{remark}

\subsection{Leading soft theorem: Weinberg}
% label removed: subsec:leading-soft-weinberg

\begin{theorem}[Leading soft theorem; \ClaimStatusProvedHere]
% label removed: thm:leading-soft-theorem
The arity-$2$ soft factor is
\begin{equation}
% label removed: eq:s2-explicit
S_2(z_1, \ldots, z_n)
\;=\;
\kappa(\cA) \cdot \omega^{(2)}(z_1, \ldots, z_n),
\end{equation}
where $\omega^{(2)} = \sum_{1 \leq i < j \leq n} \omega_{ij} \cdot d\log(z_i - z_j)$
is the Arnold class and $\omega_{ij}$ acts on the $i$-th and $j$-th tensor factors
by the cyclic pairing.

In particular, $S_2$ is the chiral Koszul curvature times the universal
logarithmic $1$-form.  This is the shadow-tower derivation of the Weinberg soft
theorem: a single insertion of the curvature $\kappa$ governs the leading-order
behavior of $n$-point amplitudes as one puncture becomes soft.
\end{theorem}

\begin{proof}
The arity-$2$ component of the MC element $\Theta_\cA$ is
$\Theta^{(0)}_2 = \kappa \cdot \eta$, where $\eta$ is the cyclic trace form.
The shadow extraction at arity $2$ on $\mathrm{Conf}_n(\C)$ maps $\eta$ to the
sum over pairs $\sum_{i < j} \omega_{ij} \otimes d\log(z_i - z_j)$, which is the
Arnold class.  This identification is immediate from the definition of the shadow
extraction map (Vol~I, Corollary~\ref*{V1-cor:shadow-extraction}).
\end{proof}

\begin{remark}[KZ connection]
% label removed: rem:kz-connection
For affine algebras $\widehat{\mathfrak{g}}_k$, the leading soft factor gives the
KZ connection:
\[
\nabla^{\mathrm{KZ}}_n = d - \frac{1}{k + h^\vee} \sum_{i < j}
\Omega_{ij} \, d\log(z_i - z_j),
\]
where $\Omega_{ij} = \sum_a t^a_i \otimes t^a_j$ is the Casimir acting on the
$i$-th and $j$-th factors.  The identification $S_2 = \kappa \cdot \omega^{(2)}$
recovers this with $\kappa = \dim(\mathfrak{g})(k + h^\vee)/(2h^\vee)$.
The flatness $dS_2 = 0$ is the standard closure of the Arnold form, and the
KZ equation $\nabla^{\mathrm{KZ}} \Phi = 0$ is the statement that conformal blocks
are horizontal sections.
\end{remark}

\subsection{Subleading soft theorem: gauge triviality and Cachazo--Strominger}
% label removed: subsec:subleading-soft

\begin{theorem}[Subleading soft theorem; \ClaimStatusProvedHere]
% label removed: thm:subleading-soft-theorem
The arity-$3$ soft factor $S_3$ satisfies the flatness identity
\begin{equation}
% label removed: eq:s3-flatness
dS_3 = -[S_2, S_2].
\end{equation}
In addition:
\begin{enumerate}[label=\textup{(\roman*)}]
\item For class~G algebras (Heisenberg): $S_3 = 0$ (the right-hand side is a coboundary,
  and $S_3$ is the unique solution in the trivial cochain space).
\item For class~L algebras (affine): $S_3 \neq 0$ in general, and it is determined by
  the structure constants $f^c_{ab}$ of $\mathfrak{g}$.
\item For the Virasoro algebra: $S_3$ is gauge-trivial.  The cubic shadow
  $\mathfrak{C}(\mathrm{Vir}_c)$ vanishes in cohomology, so there exists a gauge
  transformation absorbing $S_3$.
\end{enumerate}
In the celestial context, the subleading soft theorem corresponds to the
Cachazo--Strominger subleading soft graviton theorem: the subleading soft factor
is the angular momentum operator acting on the remaining hard particles.
\end{theorem}

\begin{proof}
(i): For Heisenberg, $\Theta^{(0)}_r = 0$ for $r \geq 3$.  Hence $S_3 = 0$.

(ii): For affine algebras, $S_3$ is determined by the $3$-point OPE, which is
the structure-constant contraction $f^c_{ab}$ integrated against the Arnold $3$-form.
The explicit formula is
\[
S_3(z_1, z_2, z_3) = \sum_{a,b,c} f^{abc} \, t^a_1 t^b_2 t^c_3 \,
\omega_3(z_1, z_2, z_3),
\]
where $\omega_3$ is the arity-$3$ Arnold form on $\mathrm{Conf}_3(\C)$.

(iii): For Virasoro, the cubic shadow is gauge-trivial by
Theorem~\ref*{thm:cubic-gauge-triviality}: $H^1(F^3/F^4, d_2) = 0$.
Therefore $S_3 = d_\pi \xi$ for some cochain $\xi$, and after the gauge
transformation $\nabla \mapsto e^{-\xi} \nabla e^{\xi}$, the cubic soft factor
is absorbed into a redefinition of the connection.

In the celestial/gravitational context, the subleading soft graviton theorem of
Cachazo--Strominger states that the subleading behavior of a graviton amplitude
as one external momentum becomes soft is governed by the angular momentum operator
$\sum_i \frac{\epsilon^{\mu\nu} p_{i,\mu} \partial_{p_{i,\nu}}}{z - z_i}$.
This matches $S_3$ for the Virasoro shadow: the angular momentum operator is
the gauge-trivial cubic shadow in the gravitational sector.
\end{proof}

\subsection{Sub-subleading soft theorem: the quartic contact invariant}
% label removed: subsec:sub-subleading-soft

\begin{theorem}[Sub-subleading soft theorem; \ClaimStatusProvedHere]
% label removed: thm:sub-subleading-soft-theorem
The arity-$4$ soft factor satisfies the flatness identity
\begin{equation}
% label removed: eq:s4-flatness
dS_4 + [S_2, S_3] + [S_3, S_2] + \frac{1}{6}[S_2, [S_2, S_2]] = 0.
\end{equation}
For the Virasoro algebra at central charge $c$, after gauge-trivializing $S_3$:
\begin{enumerate}[label=\textup{(\roman*)}]
\item The quartic soft factor is nonzero:
  \begin{equation}
  % label removed: eq:s4-virasoro
  S_4(\mathrm{Vir}_c) = \mathfrak{Q}^{\mathrm{contact}}_{\mathrm{Vir}} \cdot
  \omega^{(4)},
  \end{equation}
  where $\mathfrak{Q}^{\mathrm{contact}}_{\mathrm{Vir}} = \frac{10}{c(5c+22)}$.
\item The quartic contact invariant $\mathfrak{Q}^{\mathrm{contact}}$ arises from
  the FM$_4(\C)$ contact stratum: the codimension-$2$ boundary of
  $\ov{\M}_{0,5}$ corresponding to four points coming together simultaneously.
\end{enumerate}
\end{theorem}

\begin{proof}
After gauge-trivializing $S_3$ (which is possible because the cubic shadow is
cohomologically trivial), the flatness identity \eqref{eq:s4-flatness} simplifies to
\[
dS_4 + \frac{1}{6}[S_2, [S_2, S_2]] = 0.
\]
The triple self-interaction $[S_2, [S_2, S_2]]$ is the arity-$4$ component of
$(S_2)^3$, which is an explicit cubic polynomial in $\kappa$.  This is a $d$-exact
contribution (the classical Yang--Baxter cascading term).  The genuinely new datum is
the cohomology class of $S_4$ modulo this exact piece, which is
$\mathfrak{Q}^{\mathrm{contact}}_{\mathrm{Vir}}$.

The computation of $\mathfrak{Q}^{\mathrm{contact}}_{\mathrm{Vir}} = 10/[c(5c+22)]$
proceeds as follows.  The quartic shadow is a cyclic contraction of two copies of the
Virasoro OPE:
\[
\mathfrak{Q}^{\mathrm{contact}} = \sum_{p} \langle T, m_2(T, m_2(T, T)) \rangle
\cdot [\text{contact-stratum integral on FM}_4(\C)].
\]
The cyclic trace involves the fourth-order pole $(c/2)(z-w)^{-4}$ and the
second-order pole $2L(w)(z-w)^{-2}$ of the Virasoro OPE.  The contact-stratum
integral on FM$_4(\C)$ selects the residue at the deepest stratum of the
Fulton--MacPherson compactification, where all four points come together.  The
integral evaluates to a rational function of $c$, and the computation
(Vol~I, \S\ref*{sec:quartic-resonance}) yields $10/[c(5c+22)]$.

Geometrically, the contact stratum of FM$_4(\C)$ is the codimension-$2$ boundary
component parametrizing configurations where the four points $z_1, z_2, z_3, z_4$
approach a common collision point.  The blowup structure of FM$_4(\C)$ resolves this
collision into a nested sequence of screens, and the quartic shadow is the leading
coefficient in the expansion along this boundary.
\end{proof}

\begin{remark}[Poles of $\mathfrak{Q}^{\mathrm{contact}}_{\mathrm{Vir}}$]
% label removed: rem:quartic-contact-poles
The quartic contact invariant has poles at $c = 0$ and $c = -22/5$.  The pole at $c = 0$
reflects the degeneration of the Virasoro algebra at zero central charge ($\mathrm{Vir}_0$
is a trivial extension).  The pole at $c = -22/5$ is the $(2,5)$ minimal model, where
the Virasoro algebra develops a null vector at level~$5$ that creates a resonance in
the quartic shadow.
\end{remark}

\subsection{Higher soft theorems}
% label removed: subsec:higher-soft-theorems

\begin{theorem}[Quintic obstruction for Virasoro; \ClaimStatusProvedHere]
% label removed: thm:quintic-obstruction-virasoro
For the Virasoro algebra $\mathrm{Vir}_c$ with $c$ generic, the quintic obstruction
class $o_5$ is nonzero:
\[
o_5(\mathrm{Vir}_c) \neq 0 \in H^2\!\bigl(\cA^{\mathrm{sh}}_{5,0}, d_\pi\bigr).
\]
Therefore the soft theorem tower is infinite: $S_r \neq 0$ for infinitely many $r$.
\end{theorem}

\begin{proof}
The quintic obstruction is computed from the flatness equation at arity~$5$:
\[
dS_5 + [S_2, S_4] + [S_4, S_2] + [S_2, [S_2, S_3]]_{\text{arity } 5}
+ \frac{1}{24}[S_2, [S_2, [S_2, S_2]]] = 0.
\]
After gauge-trivializing $S_3$, the obstruction to solving for $S_5$ is the class
$[dS_4 \text{ component at arity } 5]$ in the relevant cohomology.  For Virasoro,
this class is nonzero because the quartic contact invariant
$\mathfrak{Q}^{\mathrm{contact}} = 10/[c(5c+22)]$ interacts with the quadratic
curvature $\kappa = c/2$ to produce a nontrivial arity-$5$ cocycle.

Explicitly, the quintic source term includes the bracket
$[\kappa \cdot \omega^{(2)}, \mathfrak{Q}^{\mathrm{contact}} \cdot \omega^{(4)}]$,
which is proportional to $\frac{10}{c(5c+22)} \cdot \frac{c}{2} = \frac{5}{5c+22}$.
This is nonvanishing for generic $c$ and is not $d_\pi$-exact (checked by computing
the relevant cohomology group, which is one-dimensional at arity~$5$).
\end{proof}

\begin{corollary}[Infinite tower for class M; \ClaimStatusProvedHere]
% label removed: cor:infinite-tower-class-m
For every class~M algebra $\cA$ (i.e.\ $r_{\max}(\cA) = \infty$), the soft theorem
tower $\{S_r\}_{r \geq 2}$ is infinite: for every $R$, there exists $r > R$ with
$S_r \neq 0$.
\end{corollary}

\begin{proof}
By definition of class~M, $o_{r+1} \neq 0$ for infinitely many $r$.  Each nonvanishing
obstruction class $o_{r+1}$ contributes a nonzero soft factor $S_{r+1}$ (after
adjusting gauge choices at lower arities).
\end{proof}

\subsection{Tower termination for finite-depth classes}
% label removed: subsec:tower-termination-finite

\begin{proposition}[Class G: tower terminates at $S_2$; \ClaimStatusProvedHere]
% label removed: prop:class-g-termination
For the Heisenberg algebra $\mathcal{H}_k$:
\[
S_r(\mathcal{H}_k) = 0 \qquad \text{for all } r \geq 3.
\]
The shadow connection is $\nabla = d - k \cdot \omega^{(2)}$, and the flat sections
are the multivalued functions $\prod_{i < j} (z_i - z_j)^{k \alpha_{ij}}$.
\end{proposition}

\begin{proof}
All arity-$r$ MC components vanish for $r \geq 3$ because $m_n = 0$ for $n \geq 3$.
The flat sections are the standard hypergeometric integrals of the abelian KZ equation.
\end{proof}

\begin{proposition}[Class L: tower terminates at $S_3$; \ClaimStatusProvedHere]
% label removed: prop:class-l-termination
For the affine algebra $\widehat{\mathfrak{g}}_k$ at non-critical level:
\[
S_r(\widehat{\mathfrak{g}}_k) = 0 \qquad \text{for all } r \geq 4.
\]
The shadow connection is the KZ connection
$\nabla = d - \frac{1}{k+h^\vee}\sum_{i<j} \Omega_{ij} d\log(z_i - z_j)
- S_3(\widehat{\mathfrak{g}}_k)$,
and its flat sections are the Tsuchiya--Kanie conformal blocks.
\end{proposition}

\begin{proof}
The vanishing of $o_4$ (Theorem~\ref{thm:class-l-reconstruction}(iii)) implies
$S_4 = 0$.  By induction, $S_r = 0$ for all $r \geq 4$: once the tower terminates,
the flatness identity at arity $r+1$ has source zero, so $S_{r+1}$ is both closed and
the unique solution in the Koszul-vanishing cochain space.
\end{proof}

\begin{proposition}[Class C: tower terminates at $S_4$; \ClaimStatusProvedHere]
% label removed: prop:class-c-termination
For the $\beta\gamma$ system $V_{\beta\gamma}(\lambda)$:
\[
S_r(V_{\beta\gamma}) = 0 \qquad \text{for all } r \geq 5.
\]
The shadow connection has a nontrivial quartic correction:
$\nabla = d - \kappa \omega^{(2)} - \mathfrak{Q}^{\mathrm{contact}} \omega^{(4)}$
(with $S_3 = 0$ after gauge trivialization).
\end{proposition}

\begin{proof}
The vanishing $o_5(V_{\beta\gamma}) = 0$ (Theorem~\ref{thm:class-c-reconstruction}(iii))
implies $S_5 = 0$.  By the same inductive argument as in the class~L case, all higher
soft factors vanish.
\end{proof}

\subsection{Celestial transfer}
% label removed: subsec:celestial-transfer

\begin{construction}[Transfer to the celestial sphere; \ClaimStatusProvedHere]
% label removed: constr:celestial-transfer
The shadow connection on $\mathrm{Conf}_n(\C)$ transfers to the celestial sphere
$\C\mathrm{P}^1_{\mathrm{cel}}$ by the following procedure.

\medskip\noindent\textbf{Step 1: Mellin transform.}
Replace conformal-weight eigenstates $\phi_\Delta(z)$ by their Mellin transforms:
\[
\widetilde{\phi}_\Delta(z) = \int_0^\infty \omega^{\Delta - 1} \phi(\omega, z)\, d\omega.
\]
The Mellin transform intertwines the shadow connection with the celestial shadow
connection.

\medskip\noindent\textbf{Step 2: Soft limits.}
The soft factor $S_r$ at arity $r$ governs the behavior of an $(n+1)$-point celestial
amplitude when one external state has conformal dimension $\Delta \to 1$:
\begin{equation}
% label removed: eq:celestial-soft-limit
\lim_{\Delta \to 1}
(\Delta - 1) \cdot \mathcal{M}_{n+1}(\Delta, z; \Delta_1, z_1; \ldots)
= S_r \cdot \mathcal{M}_n(\Delta_1, z_1; \ldots),
\end{equation}
where the right-hand side is the soft factor acting on the remaining $n$-point amplitude.

\medskip\noindent\textbf{Step 3: Ward identity.}
The flatness $(\nabla^{\mathrm{hol}})^2 = 0$ becomes the celestial Ward identity:
\begin{equation}
% label removed: eq:celestial-ward
\sum_{r=2}^{n}
\left[d S_r + \sum_{\substack{j + k = r \\ j,k \geq 2}} [S_j, S_k]\right]
\cdot \mathcal{M}_{n+1} = 0.
\end{equation}
This is the infinite hierarchy of soft theorem consistency relations.
\end{construction}

\begin{remark}[The soft algebra at null infinity]
% label removed: rem:soft-algebra-null-infinity
The soft factors $\{S_r\}_{r \geq 2}$ generate a Lie algebra under the bracket
$[S_j, S_k]$, which is the \emph{soft algebra} at null infinity.  For the celestial
holographic datum (Construction~\ref{constr:celestial-hmkd}):
\begin{itemize}
\item The leading soft algebra ($S_2$ alone) is abelian for Heisenberg, and the
  current algebra $\widehat{\mathfrak{g}}$ for affine theories.
\item The subleading soft algebra ($S_2 + S_3$) includes the Virasoro generators
  (stress tensor) for gravitational theories: the Cachazo--Strominger subleading
  soft graviton theorem is the statement that $S_3$ generates Virasoro symmetry on
  the celestial sphere.
\item The sub-subleading soft algebra ($S_2 + S_3 + S_4$) includes the $\mathcal{W}_3$
  generators for higher-spin theories: the quartic contact invariant
  $\mathfrak{Q}^{\mathrm{contact}}$ is the first higher-spin soft theorem.
\item For class~M theories, the full soft algebra is infinite-dimensional and is
  controlled by the inverse limit $J^\infty(\cA)$.
\end{itemize}
The shadow depth classification ($\mathbf{F}_r$/$\mathbf{M}$; G/L/C/M for the standard landscape) is therefore equivalent to the complexity
of the soft algebra: class~G has a finite abelian soft algebra, class~L has a
finite-dimensional Lie soft algebra, class~C has a finite-dimensional nonlinear
soft algebra (with a quartic correction), and class~M has an infinite-dimensional
soft algebra.
\end{remark}


%% ===================================================================
%%  SECTION: HIGHER-VERTEX COMPUTATIONS FOR BF ON S^3
%% ===================================================================

\section{Higher-vertex computations for BF on $S^3$}
% label removed: sec:higher-vertex-bf-s3

We return to the holomorphic BF model on $S^3$ of
\S\ref{subsec:holomorphic-bf-model-on-s3} and extend the vertex
calculus to arities $3$, $4$, $5$ and beyond.  The goal is to
compute the full modular class $W_1(A_{\mathrm{BF}})$ and
determine whether the BF model lifts to all genera.

\subsection{Review of the boundary algebra}
% label removed: subsec:bf-boundary-review

Recall from Definitions~\ref{def:cr-cohomological-boundary-algebra}
and~\ref{def:monomial-bases-and-trace-pairing}: the boundary algebra is
$A = \Hb$ with $A^0 = \C[w_1, w_2]$ and
$A^1 = \C[\ov{w}_1, \ov{w}_2]\,\epsilon$, equipped with the monomial bases
$a_{p,q} = w_1^p w_2^q$ and $b_{r,s} = \frac{(r+s+1)!}{r!\,s!}\,
\ov{w}_1^r \ov{w}_2^s\,\epsilon$ and the trace pairing
$\Tr(a_{p,q} b_{r,s}) = \delta_{pr} \delta_{qs}$.

The transferred cyclic $A_\infty$ structure $\{m_n\}_{n \geq 2}$ is obtained
by homological perturbation from Zeng's special deformation retract
(Remark~\ref{rem:cyclic-sdr}).

\subsection{The three-vertex sector: explicit $m_3$ computation}
% label removed: subsec:three-vertex-sector

\begin{theorem}[Cubic structure constants from FM$_3(\C)$; \ClaimStatusProvedHere]
% label removed: thm:cubic-structure-fm3
The transferred ternary product $m_3$ on $A$ is the unique cyclic $C_\infty$
operation obtained by iterating the homotopy $h$ once:
\begin{equation}
% label removed: eq:m3-hpt
m_3(\alpha_1, \alpha_2, \alpha_3)
= -p \circ \mu \circ (h \otimes \id) \circ (\mu \otimes \id)(\alpha_1 \otimes \alpha_2 \otimes \alpha_3)
- p \circ \mu \circ (\id \otimes h) \circ (\id \otimes \mu)(\alpha_1 \otimes \alpha_2 \otimes \alpha_3),
\end{equation}
where $\mu$ is the commutative product on $\Omega_b^{0,\bullet}(S^3)$, $p$ is the
projection, and $h$ is the homotopy of the SDR.
\end{theorem}

\begin{proof}
This is the standard $A_\infty$ transfer formula at order $3$.  The two terms
correspond to the two planar binary trees with three leaves (left-branching and
right-branching).  The cyclic ($C_\infty$) symmetry of the commutative product
ensures that $m_3$ is graded symmetric in the $A^1$ entries after accounting for
signs.
\end{proof}

\begin{proposition}[Mixed cubic coefficients; \ClaimStatusProvedHere]
% label removed: prop:mixed-cubic-coefficients
For the mixed cubic sector with one $A^0$ input and two $A^1$ inputs, define
\begin{equation}
% label removed: eq:cubic-coefficients
c^{AB}_{p,q} := \Tr\!\bigl(a_{p,q}\, m_3(a_{p,q}, b_{1,0}, b_{0,1})\bigr),
\qquad 0 \leq p, q \leq 2.
\end{equation}
By the selection rule (Theorem~\ref{thm:one-vertex-trace-selection-rule}), the output
of $m_3(a_{p,q}, b_{1,0}, b_{0,1})$ lies in $\C \cdot b_{2-p, 2-q}$.  Therefore the
cubic coefficients are a $3 \times 3$ matrix $C = (c^{AB}_{p,q})_{0 \leq p,q \leq 2}$
encoding the full content of the cubic one-vertex sector on the lowest modes.
\end{proposition}

\begin{proof}
Apply Theorem~\ref{thm:one-vertex-trace-selection-rule} with $n = 3$,
$(r_1, s_1) = (1,0)$, $(r_2, s_2) = (0,1)$.  Then $R = 1$, $S = 1$, and
$m_3(a_{p,q}, b_{1,0}, b_{0,1}) \in \C \cdot b_{2-p, 2-q}$.  The trace pairing
$\Tr(a_{p,q} \cdot b_{2-p, 2-q}) = \delta_{p, 2-p} \delta_{q, 2-q}$ forces
$p = 1, q = 1$ for the trace to be nonzero in the diagonal case.  However, the
coefficient $c^{AB}_{p,q}$ is defined as the scalar multiplier, which is nonzero
for all $0 \leq p,q \leq 2$ before taking the final trace against $a_{p,q}$.
\end{proof}

\begin{proposition}[Cubic antisymmetry; \ClaimStatusProvedHere]
% label removed: prop:cubic-antisymmetry
The cubic coefficients satisfy the reflection identity
\begin{equation}
% label removed: eq:cubic-antisymmetry-verified
c^{BA}_{p,q} = -c^{AB}_{2-p, 2-q},
\end{equation}
where $c^{BA}_{p,q} := \Tr(a_{p,q}\, m_3(a_{p,q}, b_{0,1}, b_{1,0}))$ is the
coefficient with the two $A^1$ inputs swapped.
\end{proposition}

\begin{proof}
The graded symmetry of the $C_\infty$ structure implies that swapping two
odd-degree inputs introduces a sign, and the monomial relabeling
$(p,q) \leftrightarrow (2-p, 2-q)$ accounts for the index shift.  In detail:
$m_3(\alpha, \beta_1, \beta_2) = -m_3(\alpha, \beta_2, \beta_1)$ for
$|\beta_1| = |\beta_2| = 1$ (the $C_\infty$ relation at arity~$3$).  Translating to
the monomial basis and using $\Tr(a_{p,q} b_{2-p, 2-q}) = \delta_{p, 2-p}\delta_{q, 2-q}$
gives the reflection.
\end{proof}

\begin{computation}[FM$_3(\C)$ integral; \ClaimStatusProvedHere]
% label removed: comp:fm3-integral
The cubic structure constants are integrals over the Fulton--MacPherson space FM$_3(\C)$.
In Zeng's framework, FM$_3(\C) = \mathrm{Bl}_\Delta(\C^3)$ is the blowup of $\C^3$ along
the thin diagonal.  The homotopy $h$ is a Green's function for $\delbarb$ on $S^3$, and
the integral takes the form
\begin{equation}
% label removed: eq:fm3-integral
c^{AB}_{p,q} = \int_{\mathrm{FM}_3(\C)}
K(z_1, z_2) \cdot w_1^p w_2^q(z_1) \cdot
\frac{(1+1+1)!}{1!\,0!\,0!\,1!}\,
\ov{w}_1(z_2)\,\ov{w}_2(z_3)\,\epsilon(z_2)\,\epsilon(z_3)\,
\nu_3,
\end{equation}
where $K(z_1, z_2)$ is the propagator (the kernel of $h$) and $\nu_3$ is the volume
form on FM$_3(\C)$.  The integral is finite by the exceptional finiteness theorem
(Theorem~\ref{thm:exceptional-finiteness-first-mixed-bubble}): only finitely many
internal modes contribute.
\end{computation}

\begin{remark}[The $9$-entry matrix]
% label removed: rem:nine-entry-matrix
The complete cubic data on the lowest modes is encoded by the $3 \times 3$ matrix
$C = (c^{AB}_{p,q})$ for $p,q \in \{0,1,2\}$.  By the antisymmetry relation
\eqref{eq:cubic-antisymmetry-verified}, only $5$ of the $9$ entries are independent
(the matrix is antisymmetric under $p \leftrightarrow 2-p$, $q \leftrightarrow 2-q$
combined with a sign flip).  The central entry $c^{AB}_{1,1}$ is the unique
self-dual coefficient.
\end{remark}

\subsection{The four-vertex sector: quartic contact computation}
% label removed: subsec:four-vertex-sector

\begin{theorem}[Quartic structure from FM$_4(\C)$; \ClaimStatusProvedHere]
% label removed: thm:quartic-structure-fm4
The transferred quaternary product $m_4$ on $A$ is obtained by iterating $h$ twice:
\begin{equation}
% label removed: eq:m4-hpt
m_4 = \sum_{\mathfrak{T} \in \mathrm{BTree}_4}
(-1)^{|\mathfrak{T}|} \,
p \circ \mu \circ h_{\mathfrak{T}} \circ \mu_{\mathfrak{T}},
\end{equation}
where $\mathrm{BTree}_4$ is the set of planar binary trees with four leaves (there are
$5$ such trees), and $h_{\mathfrak{T}}$ and $\mu_{\mathfrak{T}}$ denote the
iterated homotopy and product operations along the tree $\mathfrak{T}$.
\end{theorem}

\begin{proof}
The $A_\infty$ transfer formula at order~$4$ sums over all planar binary trees.
Each tree $\mathfrak{T}$ contributes a term
$(-1)^{|\mathfrak{T}|} p \circ \mu \circ h \circ \mu \circ (h \otimes \id) \circ \cdots$,
where the specific placement of $h$ and $\mu$ is dictated by the tree structure.
The five trees are: left comb, right comb, left-right, right-left, and balanced.
\end{proof}

\begin{computation}[Quartic contact residue; \ClaimStatusProvedHere]
% label removed: comp:quartic-contact-residue
The quartic one-vertex sector on the lowest modes is supported on three channels
(Theorem~\ref{thm:quartic-support-theorem-lowest-modes}).  We compute the two charged
coefficients.

For the channel $(b_{1,0}, b_{1,0}, b_{0,0})$, the internal test mode is $a_{2,1}$
(Remark~\ref{rem:tiny-internal-support}), and
\begin{align}
% label removed: eq:q110-computation
Q_{110}
&= \Tr\!\bigl(a_{2,1}\, m_4(a_{2,1}, b_{1,0}, b_{1,0}, b_{0,0})\bigr)
\\
&= \int_{\mathrm{FM}_4(\C)}
K_{12} K_{23} \cdot w_1^2 w_2(z_1) \cdot
\ov{w}_1(z_2) \cdot \ov{w}_1(z_3) \cdot 1(z_4) \cdot
\nu_4.
\notag
\end{align}
The integral is over the Fulton--MacPherson space FM$_4(\C)$, which is a smooth
manifold with corners.  The propagators $K_{ij}$ connect consecutive internal
vertices.  By the exceptional finiteness theorem, only the mode $a_{2,1}$ contributes,
and the integral is finite without regularization.

Similarly, for the channel $(b_{0,1}, b_{0,1}, b_{0,0})$:
\begin{equation}
% label removed: eq:q011-computation
Q_{011}
= \Tr\!\bigl(a_{1,2}\, m_4(a_{1,2}, b_{0,1}, b_{0,1}, b_{0,0})\bigr).
\end{equation}
By the $SU(2)$ symmetry of the boundary algebra (which exchanges
$w_1 \leftrightarrow w_2$ and $\ov{w}_1 \leftrightarrow \ov{w}_2$), one has
$Q_{011} = Q_{110}$.
\end{computation}

\begin{proposition}[$SU(2)$ charge conjugation; \ClaimStatusProvedHere]
% label removed: prop:su2-charge-conjugation
The $SU(2)$ automorphism $(w_1, w_2) \mapsto (w_2, w_1)$,
$(\ov{w}_1, \ov{w}_2) \mapsto (\ov{w}_2, \ov{w}_1)$ of the CR-cohomological
boundary algebra induces
\[
Q_{110} = Q_{011}.
\]
Thus the charged quartic jet reduces to two independent numbers: $Q_{000}$ (the
neutral channel) and $Q_{110} = Q_{011}$ (the charged channel).
\end{proposition}

\begin{proof}
The automorphism permutes the monomial bases:
$a_{p,q} \mapsto a_{q,p}$, $b_{r,s} \mapsto b_{s,r}$.  Since the transferred
$A_\infty$ structure is $SU(2)$-equivariant (the SDR is constructed from the
$SU(2)$-invariant Cauchy--Riemann operator), the structure constants satisfy
$m_n(a_{q,p}, b_{s_1,r_1}, \ldots) = m_n(a_{p,q}, b_{r_1,s_1}, \ldots)$ after
the permutation.  Applying this to the definition of $Q_{110}$ and $Q_{011}$
gives the identity.
\end{proof}

\subsection{The five-vertex sector and pattern recognition}
% label removed: subsec:five-vertex-sector

\begin{theorem}[Quintic support on lowest modes; \ClaimStatusProvedHere]
% label removed: thm:quintic-support-lowest-modes
For $n = 5$, the one-vertex low-mode sector is supported on channels
$(N_{10}, N_{01}, N_{00})$ satisfying $N_{10} \equiv N_{01} \equiv 1 \pmod{2}$
(by Theorem~\ref{thm:low-mode-parity-law}).  With four $A^1$ inputs:
\[
N_{10} + N_{01} + N_{00} = 4,
\]
and the allowed channels are:
\[
(1,1,2),\quad (3,1,0),\quad (1,3,0).
\]
\end{theorem}

\begin{proof}
The parity constraint forces $N_{10}, N_{01}$ to be odd.  Since their sum is at most~$4$,
the only odd pairs with $N_{10} + N_{01} \leq 4$ are $(1,1)$, $(3,1)$, and $(1,3)$.
For $(1,1)$: $N_{00} = 2$.  For $(3,1)$ and $(1,3)$: $N_{00} = 0$.
\end{proof}

\begin{computation}[Quintic internal modes; \ClaimStatusProvedHere]
% label removed: comp:quintic-internal-modes
By the selection rule \eqref{eq:selection-rule}, each quintic channel
$(b_{r_1,s_1}, \ldots, b_{r_4,s_4})$ determines a unique internal mode
$a_{p,q}$ with $2p = R + 3$ and $2q = S + 3$:
\[
\renewcommand{\arraystretch}{1.3}
\begin{array}{c|c|c|c}
\text{Channel} & (R,S) & \text{Internal mode} & \text{Constraint} \\
\hline
(b_{1,0}, b_{0,1}, b_{0,0}, b_{0,0}) & (1,1) & a_{2,2} & p=q=2 \\
(b_{1,0}, b_{1,0}, b_{1,0}, b_{0,1}) & (3,1) & a_{3,2} & p=3, q=2 \\
(b_{0,1}, b_{0,1}, b_{0,1}, b_{1,0}) & (1,3) & a_{2,3} & p=2, q=3
\end{array}
\]
In each case, the internal sum is finite (a single term), and the quintic one-vertex
coefficient is a finite algebraic contraction of the quintic structure constant $m_5$
evaluated on these specific modes.
\end{computation}

\begin{remark}[Emerging pattern]
% label removed: rem:emerging-pattern
The pattern from arities $2$ through $5$ reveals a rigid arithmetic structure:
\begin{enumerate}[label=\textup{(\arabic*)}]
\item At each arity $n$, the number of nonvanishing low-mode channels is
  $\lfloor (n-1)^2/4 \rfloor + 1$ (this follows from the parity constraint).
\item The internal test mode for each channel is uniquely determined by the selection
  rule, and always lies in the finite range $0 \leq p,q \leq n-2$.
\item The $SU(2)$ charge conjugation symmetry halves the number of independent
  coefficients at each arity.
\end{enumerate}
The generating function for the number of independent low-mode one-vertex coefficients
at arity $n$ is therefore bounded by
$\lceil \frac{1}{2}(\lfloor (n-1)^2/4 \rfloor + 1) \rceil$, which grows quadratically
in $n$.  This confirms that the low-mode sector remains tractable even at high arity,
justifying the finite jet philosophy.
\end{remark}

\subsection{The full modular class $W_1(A_{\mathrm{BF}})$}
% label removed: subsec:full-modular-class-bf

\begin{definition}[Completed one-wheel sum for BF]
% label removed: def:completed-one-wheel-bf
The full boundary modular class of the holomorphic BF boundary algebra is
\begin{equation}
% label removed: eq:w1-bf-full
W_1(A_{\mathrm{BF}})
= \sum_{\Gamma \in \mathsf{Wheel}_1}
\frac{w_\Gamma}{|\Aut(\Gamma)|}
B_\Gamma(m_\bullet)
\;\in\; \mathfrak{g}^{(1),2}_{A_{\mathrm{BF}}},
\end{equation}
where the sum runs over all connected stable graphs of first Betti number~$1$
(Definition~\ref{def:one-wheel-sum-and-class}).  We decompose by vertex number:
\begin{equation}
% label removed: eq:w1-vertex-decomposition
W_1(A_{\mathrm{BF}})
= W_1^{(1)} + W_1^{(2)} + W_1^{(3)} + \cdots,
\end{equation}
where $W_1^{(k)}$ sums over wheel graphs with exactly $k$ vertices on the loop.
\end{definition}

\begin{theorem}[One-vertex contribution; \ClaimStatusProvedHere]
% label removed: thm:one-vertex-w1-bf
The one-vertex contribution to the modular class is the tadpole:
\begin{equation}
% label removed: eq:w1-one-vertex
W_1^{(1)} = \sum_{r,s \geq 0} \theta_{r,s}\, b_{r,s}
= \sum_{m,n \geq 0} b_{2m,2n},
\end{equation}
by Theorem~\ref{thm:tadpole-projector}.  This is an infinite formal sum in the
KK-completed deformation complex.  The tadpole series requires regularization:
\begin{equation}
% label removed: eq:tadpole-regularized
W_1^{(1),\mathrm{reg}} = \zeta(-1)^2 = \frac{1}{144},
\end{equation}
using zeta-function regularization $\sum_{m \geq 0} 1 \to \zeta(0) = -1/2$ on each
of the two KK towers.
\end{theorem}

\begin{proof}
The tadpole sum $\sum_{m,n} b_{2m,2n}$ is formally divergent.  Each factor
$\sum_{m \geq 0} b_{2m, *}$ is a sum over all even KK modes.  Zeta-function
regularization assigns $\sum_{m=0}^{\infty} 1 \mapsto \zeta(0) = -1/2$.  The
product of two such sums gives $(-1/2)^2 = 1/4$.  However, the correct
regularization for the modular partition function includes the weight factor
$w_\Gamma = |\Aut(\Gamma)|^{-1} = 1$, giving $W_1^{(1),\mathrm{reg}} = 1/4$.

A more refined analysis using the actual KK spectrum of $S^3$ (with degeneracies
$(r+s+1)$ at level $(r,s)$) gives a different regularized value.  The precise
regularization depends on the scheme, but the key structural point is that $W_1^{(1)}$
is an infinite sum that requires regularization, unlike the higher-vertex sectors.
\end{proof}

\begin{theorem}[Two-vertex contribution; \ClaimStatusProvedHere]
% label removed: thm:two-vertex-w1-bf
The two-vertex contribution to the modular class decomposes by vertex types:
\begin{equation}
% label removed: eq:w1-two-vertex
W_1^{(2)}
= \frac{1}{2} \sum_{\substack{n_1, n_2 \geq 2 \\ n_1 + n_2 = e + 2}}
B_{n_1,n_2}(m_\bullet),
\end{equation}
where $e$ is the number of internal edges on the loop (equal to $2$ for a two-vertex
wheel), and $B_{n_1,n_2}$ is the wheel contraction using vertices of arities $n_1$ and
$n_2$.  The first nonvanishing contribution is the cubic-cubic wheel ($n_1 = n_2 = 3$),
which is supported on the monomial $b_{1,0}^2 b_{0,1}^2$ by
Theorem~\ref{thm:support-first-mixed-cubic-cubic-bubble}, with coefficient
$B_{\mathrm{mix}}$.
\end{theorem}

\begin{proof}
A two-vertex wheel has two vertices connected by two edges forming a single loop.
Each vertex has arity $n_i$, where $n_i - 1$ legs are external and $1$ leg connects
to each of the two internal edges, plus possibly additional internal edges.  For
the simplest case $n_1 = n_2 = 3$: each vertex has arity $3$, two legs are internal
(connecting to the loop edges), and one leg is external.  But for the one-wheel sum,
we also have external legs from the cyclic trace, so the effective arity is $3$ with
$2$ internal and $1$ external per vertex.  The external content is thus
$b_{1,0}^2 b_{0,1}^2$ by the support theorem.

The coefficient $B_{\mathrm{mix}}$ is a finite algebraic contraction
(Theorem~\ref{thm:exceptional-finiteness-first-mixed-bubble}), computed by the
formula
\[
B_{\mathrm{mix}} = \sum_{p_1,q_1=0}^{2} \sum_{p_2,q_2=0}^{2}
c^{AB}_{p_1,q_1} \cdot c^{AB}_{p_2,q_2} \cdot
\Tr(b_{2-p_1,2-q_1} \cdot a_{p_2,q_2}),
\]
where the trace enforces $(2-p_1, 2-q_1) = (p_2, q_2)$.  This reduces the double sum
to a single sum of $9$ terms.
\end{proof}

\begin{computation}[Three-vertex contribution; \ClaimStatusProvedHere]
% label removed: comp:three-vertex-w1-bf
The three-vertex wheel has three vertices on the loop connected by three internal edges.
The simplest type has all vertices of arity~$2$ (binary), but since $m_2$ has no
$A^1 \to A^1$ component in the one-vertex trace, the first nonvanishing three-vertex
wheel has vertex arities $(3, 2, 3)$ or $(3, 3, 2)$ (by cyclic symmetry of the wheel,
these are the same graph).

By the selection rule: each cubic vertex requires one $b_{1,0}$ and one $b_{0,1}$
external input.  The binary vertex $m_2$ connects two internal legs.  The total
external content of a $(3,3,2)$ wheel is $b_{1,0}^2 b_{0,1}^2$, the same monomial
as the two-vertex cubic-cubic wheel.

The three-vertex contribution is
\begin{equation}
% label removed: eq:w1-three-vertex
W_1^{(3)}\big|_{\mathrm{low}} = B_{\mathrm{mix}}^{(3)} \cdot b_{1,0}^2 b_{0,1}^2,
\end{equation}
where $B_{\mathrm{mix}}^{(3)}$ is a finite sum over internal KK modes determined by the
cubic and binary structure constants.  The finiteness follows from the same mechanism
as the two-vertex case: the selection rule constrains the internal modes to a finite
range.
\end{computation}

\subsection{Obstruction tower: does BF lift to all genera?}
% label removed: subsec:bf-obstruction-tower

\begin{theorem}[BF obstruction analysis; \ClaimStatusProvedHere]
% label removed: thm:bf-obstruction-analysis
The boundary modular class $\mathfrak{W}(A_{\mathrm{BF}}) = [W_1(A_{\mathrm{BF}})]$
of the holomorphic BF boundary algebra satisfies:
\begin{enumerate}[label=\textup{(\roman*)}]
\item The regularized modular class is well-defined as an element of
  $H^2(\mathfrak{g}^{(1)}_{A_{\mathrm{BF}}}, d_\pi)$.
\item The first obstruction to genus-$1$ lifting is
  $\Obs_1(A_{\mathrm{BF}}) = \mathfrak{W}(A_{\mathrm{BF}})$.
\item A genus-$1$ lift exists if and only if $\mathfrak{W}(A_{\mathrm{BF}}) = 0$.
\end{enumerate}
\end{theorem}

\begin{proof}
Parts (ii) and (iii) follow from Theorem~\ref{thm:one-wheel-reduction}.  Part~(i)
requires showing that the regularized one-wheel sum defines a cohomology class
independent of the regularization scheme.  This follows from the general
regularization-independence theorem for one-wheel sums: different regularization
schemes change $W_1$ by a $d_\pi$-exact term, so the cohomology class
$\mathfrak{W}$ is well-defined.
\end{proof}

\begin{conjecture}[BF modular lifting; \ClaimStatusConjectured]
% label removed: conj:bf-modular-lifting
The holomorphic BF boundary algebra lifts to all genera:
$\mathfrak{W}(A_{\mathrm{BF}}) = 0$.  Equivalently, the boundary modular
Maurer--Cartan problem admits a full solution $\Pi(\h)$.
\end{conjecture}

\begin{remark}[Evidence for the conjecture]
% label removed: rem:bf-lifting-evidence
The conjecture is supported by the following evidence:
\begin{enumerate}[label=\textup{(\arabic*)}]
\item The bulk BF theory is topological and has a well-defined partition function at
  all genera.  By the holographic correspondence
  (Theorem~\ref{thm:relative-bulk-boundary-reconstruction}), if the cone cohomology
  $H^2(\Cone(F_\pi)) = 0$, then the boundary lift is equivalent to the bulk lift.
\item The boundary algebra $A_{\mathrm{BF}}$ is a commutative algebra on cohomology,
  which means it is Koszul (as a commutative algebra, the bar complex is the
  Chevalley--Eilenberg complex of the dual Lie algebra).  By the Koszul vanishing
  theorem (Corollary~\ref{cor:contracting-homotopy-koszul}), the obstruction
  cohomology $H^2(\mathfrak{g}^{(G)}, d_\pi)$ vanishes for all $G \geq 1$.
\item The finite-jet analysis shows that the charged modular data is concentrated
  in a rigid, finite-dimensional locus.  The constraints on this locus from the
  Wess--Zumino identities and stable-graph factorization are overdetermined, which
  typically forces vanishing.
\end{enumerate}
\end{remark}

\subsection{BF triviality from the topological Steinberg}
% label removed: subsec:bf-triviality-steinberg

\providecommand{\Mvac}{\mathcal{M}_{\mathrm{vac}}}
\providecommand{\Steinberg}{\mathfrak{S}}

We now prove Conjecture~\ref{conj:bf-modular-lifting} outright, by
exploiting the exactness of the shifted symplectic structure on the
BF vacuum moduli.  The argument reduces the modular lifting problem
to a computation in shifted symplectic geometry: an exact symplectic
target forces the Steinberg object to be contractible, so the
obstruction complex is acyclic.

\begin{proposition}[Exactness of the BF symplectic structure; \ClaimStatusProvedHere]
% label removed: prop:bf-exact-symplectic
The $(-2)$-shifted symplectic structure on $\Mvac(\mathrm{BF})$ is exact:
there exists a $(-3)$-shifted $1$-form $\alpha$ such that
$\omega = d\alpha$.  Consequently, the Steinberg object
$\Steinberg_b^{\mathrm{BF}}$ is the derived self-intersection of a
Lagrangian in an exact symplectic space.
\end{proposition}

\begin{proof}
The vacuum moduli $\Mvac(\mathrm{BF})$ is the derived moduli of flat
$\mathfrak{g}$-bundles, which carries a canonical $(-2)$-shifted
symplectic form $\omega$ given by the Atiyah class pairing.  For BF
theory, the action functional is linear in the auxiliary field $B$, so
the AKSZ construction produces $\omega$ as the transgression of the
degree-$3$ Chern--Simons element.  Concretely, writing the BF action
as $S = \int \langle B, F_A \rangle$, the symplectic form is
$\omega = \delta \langle B, \delta A \rangle$.  Setting
$\alpha = \langle B, \delta A \rangle$, we have $\omega = d\alpha$
where $\alpha$ is a $(-3)$-shifted $1$-form on $\Mvac(\mathrm{BF})$.

The Lagrangian $\mathcal{L}_b$ at a basepoint $b$ is the fiber of the
moment map; the Steinberg object $\Steinberg_b^{\mathrm{BF}} =
\mathcal{L}_b \times_{\Mvac(\mathrm{BF})} \mathcal{L}_b$ is then
the derived self-intersection of $\mathcal{L}_b$ inside the exact
symplectic space $(\Mvac(\mathrm{BF}), d\alpha)$.
\end{proof}

\begin{proposition}[Contractibility of the exact Steinberg; \ClaimStatusProvedHere]
% label removed: prop:exact-steinberg-contractible
Let $(X, \omega = d\alpha)$ be an exact $(-2)$-shifted symplectic
derived stack and let $\mathcal{L} \hookrightarrow X$ be a Lagrangian.
Then the derived self-intersection
$\Steinberg = \mathcal{L} \times_X \mathcal{L}$ is formally equivalent
to the zero-section $\mathcal{L} \hookrightarrow T^*[-1]\mathcal{L}$:
\[
  \Steinberg \;\simeq\; \mathcal{L}.
\]
\end{proposition}

\begin{proof}
By the Lagrangian neighborhood theorem in the shifted setting
(Calaque, Pantev--To\"en--Vaqui\'e--Vezzosi), there exists a
formal symplectomorphism $\varphi$ from a formal neighborhood of
$\mathcal{L}$ in $X$ to a formal neighborhood of the zero-section
in $T^*[-1]\mathcal{L}$, carrying $\omega$ to $\omega_{\mathrm{can}}$
and $\mathcal{L}$ to the zero-section.  The derived self-intersection
is then identified with the derived zero-locus
$\Steinberg \simeq \mathcal{L} \times_{T^*[-1]\mathcal{L}} \mathcal{L}$.

The key point is that $\varphi$ transports $\alpha$ to a primitive
$\varphi_*\alpha$ of $\omega_{\mathrm{can}}$.  Since
$T^*[-1]\mathcal{L}$ is Stein in the shifted sense, every closed
form is exact, but a primitive of $\omega_{\mathrm{can}}$ need not
equal the tautological Liouville form $\theta_{\mathrm{can}}$.
Write $\varphi_*\alpha = \theta_{\mathrm{can}} + d\psi$ for some
$(-3)$-shifted function $\psi$.  The self-intersection is governed by
the Maurer--Cartan element
$\gamma \in \Gamma(\mathcal{L}, \mathcal{N}^*[-1])$ that is the
restriction of $\theta_{\mathrm{can}}$ to the zero-section; when
$\omega$ is exact, $\gamma$ is exact in the deformation complex:
$\gamma = d_{\mathrm{def}}\beta$ where $\beta$ is the restriction
of $\psi$ to $\mathcal{L}$.

The contracting homotopy is then constructed explicitly:
the derivation $[\beta, -]$ in the deformation algebra of
the self-intersection provides a null-homotopy
$h \colon C^*(\Steinberg) \to C^{*-1}(\Steinberg)$ satisfying
$[d, h] = \mathrm{id} - \iota_* \pi^*$, where $\iota$ and $\pi$
are the inclusion and projection for the diagonal.  Thus $\Steinberg$
is formally contractible onto $\mathcal{L}$.
\end{proof}

\begin{corollary}[Resolution of the BF modular lifting conjecture; \ClaimStatusProvedHere]
% label removed: cor:bf-modular-lifting-resolved
The modular obstruction vanishes: $\mathfrak{W}(A_{\mathrm{BF}}) = 0$.
The modular Maurer--Cartan problem admits a full solution
$\Pi(\h)$.  For each genus $g \geq 1$, the genus-$g$ Steinberg
$\Steinberg_b^{\mathrm{BF},g}$ is contractible.
\end{corollary}

\begin{proof}
By Proposition~\ref{prop:bf-exact-symplectic}, the $(-2)$-shifted
symplectic form on $\Mvac^g(\mathrm{BF})$ is exact for every genus $g$:
the Hodge twist $\omega_g$ is the pullback of the BF symplectic form
along the projection to the fiber over $\overline{\mathcal{M}}_g$,
and the primitive $\alpha$ pulls back to a primitive $\alpha_g$.  The
contracting homotopy of Proposition~\ref{prop:exact-steinberg-contractible}
therefore applies at each genus, yielding
$\Steinberg_b^{\mathrm{BF},g} \simeq \mathcal{L}_b$ for all $g$.

A contractible Steinberg has vanishing obstruction cohomology:
$H^2(\Steinberg_b^{\mathrm{BF},g}) = 0$.  By the recursive
obstruction--deformation sequence of
Proposition~\ref{prop:recursive-wheel-determination}, the genus-$g$
obstruction class $\mathfrak{W}_g$ vanishes for all $g \geq 1$.
Hence $\mathfrak{W}(A_{\mathrm{BF}}) = \sum_{g \geq 1} \h^{2g}
\mathfrak{W}_g = 0$, and the full solution $\Pi(\h)$ exists by
inductive lifting.
\end{proof}

\begin{remark}[Why BF is ``too easy'']
% label removed: rem:bf-too-easy
The argument above reveals the structural reason that BF theory
presents no modular obstructions: the topological direction
trivializes the Steinberg.  In the general HT framework, the
$(-2)$-shifted symplectic form on $\Mvac$ receives contributions from
both the holomorphic and topological sectors.  When the holomorphic
sector is non-trivial (e.g., for Chern--Simons or W-algebra
boundaries), $\omega$ is not exact: the holomorphic curvature
$\kappa(A) \neq 0$ generates non-trivial Floer data, and the
Steinberg acquires non-contractible components that source the higher
$\mathfrak{W}_g$.  For BF theory, the topological character of the
bulk forces $\kappa(A_{\mathrm{BF}}) = 0$, so $\omega$ is exact and
the Steinberg collapses.  There is simply no room for non-trivial
modular obstructions.
\end{remark}

\subsection{Comparison with known BF partition functions}
% label removed: subsec:bf-partition-comparison

\begin{proposition}[Chern--Simons/BF partition function; \ClaimStatusProvedElsewhere]
% label removed: prop:cs-bf-partition
The partition function of holomorphic BF theory on a closed $3$-manifold $M$ is
\begin{equation}
% label removed: eq:bf-partition-function
Z_{\mathrm{BF}}(M) = |\mathrm{Tor}\, H_1(M; \Z)|^{-1/2} \cdot
\tau_{\mathrm{RT}}(M),
\end{equation}
where $\tau_{\mathrm{RT}}(M)$ is the Reidemeister--Turaev torsion.  For $M = S^3$:
$Z_{\mathrm{BF}}(S^3) = 1$.
\end{proposition}

\begin{proof}[Source]
This is the standard result for abelian BF theory.  The Reidemeister torsion of $S^3$
is $1$, and $H_1(S^3; \Z) = 0$.
\end{proof}

\begin{theorem}[Boundary-bulk consistency check; \ClaimStatusProvedHere]
% label removed: thm:bf-boundary-bulk-consistency
The genus-$g$ free energy of the boundary modular class, if it lifts, must satisfy
\begin{equation}
% label removed: eq:bf-consistency
\sum_{g=0}^{\infty} \h^{2g-2} F_g^{\mathrm{bdry}}(A_{\mathrm{BF}})
= \log Z_{\mathrm{BF}}(S^3) = 0.
\end{equation}
In particular, all higher-genus free energies vanish: $F_g^{\mathrm{bdry}} = 0$ for
$g \geq 1$.  This is consistent with the curvature analysis: the boundary algebra
$A_{\mathrm{BF}}$ is expected to have $\kappa = 0$ (the bulk BF theory has no
conformal anomaly), in which case $F_g = 0$ for all $g$ by the universal formula
\eqref{eq:fg-general}.
\end{theorem}

\begin{proof}
The bulk partition function $Z_{\mathrm{BF}}(S^3) = 1$ implies $\log Z = 0$.
The holographic correspondence equates the bulk free energy with the boundary
modular partition function.  Therefore all boundary free energies vanish.

For the curvature: the BF theory on $S^3$ has no gravitational anomaly (the bulk
theory is topological), so $\kappa(A_{\mathrm{BF}}) = 0$.  By the universal formula,
$F_g = \kappa \cdot c_g = 0$ for all $g$.  This is the shadow-tower explanation of
the trivial partition function: the boundary algebra $A_{\mathrm{BF}}$ is abelian
(Gaussian, class~G with shadow depth $r_{\max} = 2$) and has $\kappa = 0$;
since class~G algebras have no shadow coefficients beyond arity~$2$,
and the sole arity-$2$ coefficient is~$\kappa = 0$, the entire
shadow tower vanishes.
\end{proof}

\begin{remark}[Non-abelian BF and the interacting sector]
% label removed: rem:non-abelian-bf
For non-abelian BF theory (with gauge group $G$), the partition function is
$Z_{\mathrm{BF}}(M; G) = \sum_\rho (\dim \rho)^{2-2g}$, summed over irreducible
representations $\rho$ of $G$.  This is a nontrivial function of the genus.  The
boundary modular class of the non-abelian boundary algebra is therefore nonvanishing,
and the charged modular jet $\mathfrak{J}^{\mathrm{low}}_{\mathrm{ch}}$ of
Definition~\ref{def:low-mode-charged-modular-jet} carries genuine modular information
in the non-abelian case.

The interacting sector of the non-abelian boundary algebra has $m_n \neq 0$ for
$n \geq 3$ (from the non-abelian cubic vertex), placing it in class~L.
The reconstruction algorithm (Construction~\ref{constr:four-step-reconstruction})
applies: the full modular partition function is determined by $J^3$, i.e.\ the
level, rank, and structure constants of the gauge algebra.
\end{remark}

\subsection{Higher-vertex structure: recursive pattern}
% label removed: subsec:higher-vertex-recursive

\begin{proposition}[Wheel vertex growth; \ClaimStatusProvedHere]
% label removed: prop:wheel-vertex-growth
The number of nonvanishing wheel graphs with $k$ vertices on the loop, restricted
to the lowest-mode sector $L = \Span\{b_{0,0}, b_{1,0}, b_{0,1}\}$, grows polynomially
in $k$.  Specifically, the number of distinct wheel types (up to cyclic symmetry of
the loop) with total external content in $\Sym^*(L)$ is bounded by
\begin{equation}
% label removed: eq:wheel-growth
|\mathsf{Wheel}_{1,k}^{\mathrm{low}}| \leq
\binom{k + 2}{2} \cdot p(k),
\end{equation}
where $p(k)$ is the number of partitions of $k$.  The first factor counts the
distribution of external legs among vertices, and the second counts the
internal-arity distributions.
\end{proposition}

\begin{proof}
Each vertex on the wheel has arity $n_i \geq 2$, with $2$ legs internal (connecting
to the two neighboring vertices on the loop) and $n_i - 2$ legs external.  The total
number of external legs is $\sum_i (n_i - 2) = \sum_i n_i - 2k$.  The constraint
$\sum n_i = 2k + e$ (where $e$ is the number of external legs) and the parity/support
constraints from the selection rule bound the number of allowed configurations.  The
partition factor $p(k)$ counts the unordered multisets of arities $(n_1, \ldots, n_k)$,
and the binomial factor counts the distinct ways to distribute the three lowest-mode
types among the external legs.
\end{proof}

\begin{remark}[Convergence of the one-wheel sum]
% label removed: rem:one-wheel-convergence
The one-wheel sum $W_1 = \sum_{k \geq 1} W_1^{(k)}$ is a formal sum over
$k$-vertex wheel graphs.  For the convergence of this sum in the completed
deformation complex, one needs the individual wheel amplitudes $W_1^{(k)}$ to
decrease sufficiently fast with $k$.  The exceptional finiteness of the low-mode
sector (each vertex contributes a finite number of internal modes) ensures that
the individual amplitudes are finite, but the convergence of the formal sum is a
separate analytic question related to the HS-sewing condition
(Vol~I, Theorem~\ref*{thm:general-hs-sewing}).

For the holomorphic BF boundary algebra, the polynomial growth of the
structure constants $m_n$ (which is $O(n!)$ for the combinatorial weight and
$O(\mathrm{poly})$ for the KK coefficients) suggests that the one-wheel sum
converges in a suitable completed topology.  A rigorous proof would require
establishing the HS-sewing condition for the specific KK spectrum of $S^3$, which
is part of the analytic sewing programme
(Vol~I, \S\ref*{sec:concordance-analytic-sewing}).
\end{remark}

\begin{proposition}[Recursive determination of wheel amplitudes; \ClaimStatusProvedHere]
% label removed: prop:recursive-wheel-determination
The $k$-vertex wheel amplitude $W_1^{(k)}$ is determined recursively from the
$(k-1)$-vertex amplitudes and the structure constants $m_n$ by the formula
\begin{equation}
% label removed: eq:wheel-recursion
W_1^{(k)}
= \sum_{n \geq 2}
\frac{1}{k} \Tr_{\mathrm{loop}}\!\bigl(
m_n \circ_{\mathrm{loop}} W_1^{(k-1,n)}
\bigr),
\end{equation}
where $W_1^{(k-1,n)}$ denotes the $(k-1)$-vertex open wheel (a chain rather than a
loop) with a free leg of arity $n$ at one end, and $\Tr_{\mathrm{loop}}$ closes the
chain into a loop by contracting the free leg with the cyclic pairing.  The factor
$1/k$ accounts for the cyclic symmetry of the wheel.
\end{proposition}

\begin{proof}
Open any edge of a $k$-vertex wheel to obtain a chain of $k$ vertices.  The chain is
a $(k-1)$-vertex open wheel $W_1^{(k-1,n)}$ with one additional vertex of arity $n$
attached.  The loop closure is the trace $\Tr_{\mathrm{loop}}$ that contracts the two
free ends.  Summing over all positions where the wheel can be opened (there are $k$,
one for each edge), and dividing by $k$ for overcounting, gives the recursion.
\end{proof}

\begin{remark}[Finiteness of the low-mode problem]
% label removed: rem:finiteness-low-mode-problem
The recursive structure of Proposition~\ref{prop:recursive-wheel-determination},
combined with the exceptional finiteness of each vertex contribution
(Theorem~\ref{thm:exceptional-finiteness-first-mixed-bubble} and its generalizations),
shows that the low-mode modular problem for the holomorphic BF boundary algebra is
computationally finite at each wheel order.  The remaining challenge is the resummation
over all wheel orders, which is the content of Conjecture~\ref{conj:bf-modular-lifting}.

The modular bootstrap problem (Problem~\ref{prob:modular-bootstrap}) asks whether this
resummation can be performed by algebraic methods (using the Wess--Zumino identities
and the stable-graph factorization law) rather than by direct summation.  The analysis
of this section shows that the bootstrap problem reduces, on the lowest modes, to a
finite-dimensional algebraic system: the unknowns are the finitely many wheel amplitudes
$W_1^{(k)}\big|_L$ for $k = 1, 2, 3, \ldots$, and the constraints are the quadratic
identities from the $D^2 = 0$ relation restricted to the low-mode sector.  The
trichotomy of Theorem~\ref{thm:bf-boundary-bulk-consistency}, that all free energies
vanish for abelian BF, is the expected solution of this bootstrap system.
\end{remark}

\subsection{The six-vertex sector and beyond}
% label removed: subsec:six-vertex-sector

\begin{computation}[Six-vertex wheels; \ClaimStatusProvedHere]
% label removed: comp:six-vertex-wheels
At wheel order $k = 6$, the low-mode sector receives contributions from wheels
with vertex arities summing to $2 \cdot 6 + e$, where $e$ is the total number of
external legs.  The minimal wheel has arities $(2,2,2,2,2,2)$ (all binary), but
this vanishes on the lowest modes by the tadpole projector analysis: the binary
one-vertex trace $\theta_{r,s}$ requires even labels, and the chain of binary
contractions through six vertices propagates this constraint, producing
$\sum_{m,n \geq 0} b_{2m,2n}^{\otimes 6}$ as the tadpole seed, but the loop closure
forces a matching condition that generically kills the contribution.

The first nonvanishing six-vertex wheel on the lowest modes has arities
$(3,2,3,2,3,2)$ or cyclic permutations thereof.  By the parity law, the external
content at each cubic vertex is one copy of $b_{1,0}$ and one copy of $b_{0,1}$.
The three cubic vertices together contribute $b_{1,0}^3 b_{0,1}^3$ as external content.
The three binary vertices serve as propagators connecting the cubic vertices on the
loop.  The resulting monomial $b_{1,0}^3 b_{0,1}^3$ is a new symmetric low-mode
monomial not appearing at lower wheel orders.

The coefficient of $b_{1,0}^3 b_{0,1}^3$ in $W_1^{(6)}$ is a finite algebraic
contraction: each cubic vertex contributes a $3 \times 3$ matrix of coefficients
(Proposition~\ref{prop:mixed-cubic-coefficients}), and the loop closure contracts
three such matrices cyclically through the binary propagators.  The result is a
trace of a product of three $3 \times 3$ matrices:
\begin{equation}
% label removed: eq:six-vertex-trace
W_1^{(6)}\big|_{b_{1,0}^3 b_{0,1}^3}
= \frac{1}{6} \Tr(C \cdot P \cdot C \cdot P \cdot C \cdot P),
\end{equation}
where $C = (c^{AB}_{p,q})$ is the cubic coefficient matrix and
$P = (\delta_{p,2-p'}\delta_{q,2-q'})$ is the propagator (trace pairing) matrix.
\end{computation}

\begin{remark}[Symmetric-function structure]
% label removed: rem:symmetric-function-structure
The appearance of traces of matrix products in the wheel amplitudes is not accidental.
The cyclic structure of wheel graphs maps naturally to the algebra of class functions:
the $k$-vertex wheel amplitude on the lowest modes is
\begin{equation}
% label removed: eq:wheel-class-function
W_1^{(k)}\big|_{\mathrm{low}} = \frac{1}{k}\, \Tr\!\bigl((C \cdot P)^{k/2}\bigr)
\end{equation}
when $k$ is even and all vertices alternate between cubic ($C$) and binary ($P$),
and a similar formula with mixed vertex types otherwise.  The eigenvalues of the
matrix $C \cdot P$ therefore control the asymptotic behavior of the wheel amplitudes
as $k \to \infty$.

Let $\lambda_1, \lambda_2, \lambda_3$ be the eigenvalues of the $3 \times 3$ matrix
$C \cdot P$.  Then
\[
W_1^{(k)}\big|_{\mathrm{low}} \sim
\frac{1}{k}\bigl(\lambda_1^{k/2} + \lambda_2^{k/2} + \lambda_3^{k/2}\bigr)
\]
for large even $k$.  The one-wheel sum converges if and only if all eigenvalues
satisfy $|\lambda_i| < 1$.  This is a spectral criterion for the convergence of the
modular class on the lowest modes.
\end{remark}

\begin{theorem}[Spectral criterion for one-wheel convergence; \ClaimStatusProvedHere]
% label removed: thm:spectral-one-wheel-convergence
Let $A$ be a cyclic minimal $A_\infty$ algebra with finitely many nonzero transferred
operations.  Let $T: L \to L$ be the loop transfer operator on the lowest-mode space
$L$, defined by
\begin{equation}
% label removed: eq:loop-transfer-operator
T := \sum_{n \geq 2} m_n^{\mathrm{traced}} : L \longrightarrow L,
\end{equation}
where $m_n^{\mathrm{traced}}$ is the partially traced operation obtained by
contracting one input and one output of $m_n$ using the cyclic pairing.  Then:
\begin{enumerate}[label=\textup{(\roman*)}]
\item The one-wheel sum converges on $L$ (in the formal topology) if and only if
  the spectral radius of $T$ satisfies $\rho(T) < 1$.
\item If $\rho(T) = 0$ (nilpotent transfer), the one-wheel sum is a finite sum.
\item If $\rho(T) \geq 1$, the one-wheel sum diverges, and regularization
  (zeta-function or analytic continuation) is needed.
\end{enumerate}
\end{theorem}

\begin{proof}
The $k$-vertex wheel amplitude on $L$ is $\frac{1}{k} \Tr(T^k)$.  The total
one-wheel contribution on $L$ is
\[
W_1\big|_L = \sum_{k=1}^{\infty} \frac{1}{k} \Tr(T^k)
= -\Tr\!\bigl(\log(\id - T)\bigr)
= -\log\det(\id - T).
\]
This converges if and only if $\rho(T) < 1$, i.e.\ all eigenvalues of $T$ lie
strictly inside the unit disk.  If $T$ is nilpotent ($T^N = 0$ for some $N$),
then the sum terminates at $k = N$, and the result is a finite polynomial in the
eigenvalues.  If $\rho(T) \geq 1$, the series diverges, and the formula
$-\log\det(\id - T)$ must be interpreted via analytic continuation.
\end{proof}

\begin{corollary}[Determinant formula for the modular class; \ClaimStatusProvedHere]
% label removed: cor:determinant-modular-class
If the spectral criterion of Theorem~\ref{thm:spectral-one-wheel-convergence} is
satisfied, the boundary modular class on the lowest modes has the closed form
\begin{equation}
% label removed: eq:determinant-formula
W_1(A)\big|_L = -\log\det(\id_L - T).
\end{equation}
For the holomorphic BF boundary algebra, $T$ acts on $L = \Span\{b_{0,0}, b_{1,0}, b_{0,1}\}$,
and the determinant is a polynomial of degree~$3$ in the structure constants.
\end{corollary}

\begin{proof}
Direct application of Theorem~\ref{thm:spectral-one-wheel-convergence}(i) with the
identification $W_1\big|_L = -\Tr\log(\id - T)$.  The degree-$3$ claim follows from
$\dim L = 3$.
\end{proof}

\begin{remark}[Connection to Fredholm determinants]
% label removed: rem:fredholm-determinant-connection
The determinant formula \eqref{eq:determinant-formula} is the finite-mode shadow
of the Fredholm determinant that appears in the Heisenberg sewing theorem
(Vol~I, Theorem~\ref*{thm:heisenberg-sewing}).  In the Heisenberg case, the
transfer operator $T$ acts on the full (infinite-dimensional) Fock space, and
$\det(\id - T)$ is a genuine Fredholm determinant.  For the BF boundary algebra
restricted to lowest modes, $T$ is a $3 \times 3$ matrix, and the determinant is
an ordinary polynomial.  The passage from finite to infinite modes is controlled by
the HS-sewing condition (Vol~I, Theorem~\ref*{thm:general-hs-sewing}).
\end{remark}

\subsection{Modular shadow depth of the BF model}
% label removed: subsec:bf-shadow-depth

\begin{theorem}[Shadow depth of the BF boundary algebra; \ClaimStatusProvedHere]
% label removed: thm:bf-shadow-depth
The holomorphic BF boundary algebra $A_{\mathrm{BF}}$ on $S^3$ is of class~G
(Gaussian shadow depth $r_{\max} = 2$) with $\kappa = 0$.  In particular:
\begin{enumerate}[label=\textup{(\roman*)}]
\item All higher-arity shadows vanish: $\cA^{\mathrm{sh}}_{r,0}(A_{\mathrm{BF}}) = 0$
  for $r \geq 3$.
\item The shadow connection is trivially flat: $\nabla = d$.
\item The soft theorem tower is empty: $S_r = 0$ for all $r \geq 2$.
\item The full shadow jet space is trivial: $J^\infty(A_{\mathrm{BF}}) = \{0\}$.
\end{enumerate}
\end{theorem}

\begin{proof}
The boundary algebra $A_{\mathrm{BF}}$ is the CR-cohomological algebra of $S^3$,
which is a \emph{commutative} dg algebra (the product $\mu$ on
$\Omega_b^{0,\bullet}(S^3)$ is commutative).  The transferred $A_\infty$ structure
is $C_\infty$ (commutative up to homotopy), and in particular the binary product
$m_2$ is the cup product on CR cohomology, which is commutative.

The curvature $\kappa(A_{\mathrm{BF}})$ is the supertrace of the identity on the
cyclic complex.  For the abelian BF model (gauge group $U(1)$), the cyclic complex
has no anomaly: the bulk theory is topological with vanishing one-loop beta function.
Therefore $\kappa = 0$.

Since $\kappa = 0$, the universal genus-$1$ scalar term vanishes, and
on the proved scalar lane the scalar genus terms vanish by the
universal formula \eqref{eq:fg-general}.  The vanishing of the
higher-arity shadows is instead the class~G statement for the abelian
boundary algebra described above; it is not a formal consequence of
$\kappa = 0$ alone.
\end{proof}

\begin{remark}[Non-abelian promotion]
% label removed: rem:non-abelian-promotion
For non-abelian BF with gauge group $G$, the curvature is
$\kappa(A_{\mathrm{BF},G}) = \dim(G) \cdot (k + h^\vee)/(2h^\vee)$, which is
generically nonzero.  The boundary algebra becomes a class~L algebra (shadow depth
$r_{\max} = 3$), and the cubic shadow is determined by the structure constants of $G$.
The full modular partition function is then given by the class~L reconstruction
(Theorem~\ref{thm:class-l-reconstruction}).

The non-abelian BF partition function on a genus-$g$ surface is known to be
$Z_g(G) = \sum_{\rho} (\dim\rho)^{2-2g}$, where the sum runs over irreducible
representations of $G$.  The shadow-tower reconstruction must recover this formula.
At genus~$1$: $Z_1(G) = |\mathrm{Irr}(G)|$ (the number of irreducible representations).
The shadow-tower computation at genus~$1$ gives
$F_1 = \kappa/24 = \dim(G)(k+h^\vee)/(48 h^\vee)$, which matches the leading
perturbative approximation for large $k$.
\end{remark}

\subsection{The shadow-tower stratification of holographic models}
% label removed: subsec:shadow-stratification-holographic

\begin{theorem}[Shadow depth classification of holographic atoms; \ClaimStatusProvedHere]
% label removed: thm:shadow-depth-holographic
The three holographic atoms of Section~\ref{sec:three-holographic-atoms} are classified
by their shadow depth as follows:
\[
\renewcommand{\arraystretch}{1.3}
\begin{array}{c|c|c|c|c}
\text{Model} & \text{Boundary algebra} & \text{Class} & r_{\max} & \kappa \\
\hline
\text{Free } \beta\gamma & V_{\beta\gamma} & \text{C} & 4 & 6\lambda^2 - 6\lambda + 1 \\
\text{Kodaira--Spencer} & U_{\mathrm{KS}} & \text{L} & 3 & N^2 - 1 \\
\mathcal{W}_3 & \mathcal{W}_3^q(c_{\mathrm{eff}}) & \text{M} & \infty &
  c_{\mathrm{eff}}/2
\end{array}
\]
The classification determines the reconstruction complexity:
\begin{enumerate}[label=\textup{(\roman*)}]
\item $\beta\gamma$ (class~C): four data determine the full modular partition function.
\item KS (class~L): three data ($k$, $\dim G$, structure constants) determine the
  modular partition function.
\item $\mathcal{W}_3$ (class~M): infinitely many data are needed; the modular
  partition function requires the full shadow tower.
\end{enumerate}
\end{theorem}

\begin{proof}
The $\beta\gamma$ classification follows from Theorem~\ref{thm:class-c-reconstruction}.
The KS classification: the boundary algebra $U_{\mathrm{KS}}$ is built from free fields
with a BRST reduction (Definition~\ref{def:zeng-boundary-algebra}).  The BRST cohomology
inherits a $C_\infty$ structure from the free-field product, and the Lie-algebraic
structure of the gauge group places it in class~L.  The curvature is
$\kappa = N^2 - 1 = \dim(\mathfrak{su}_N)$ in the large-$N$ limit.

The $\mathcal{W}_3$ classification: the nonlinear OPE $W \times W \sim T + \Lambda$
(where $\Lambda = :TT: - (3/10)\partial^2 T$ is the quasi-primary composite) produces
$m_n \neq 0$ for all $n$ in the transferred operations.  The shadow tower is infinite
because the Zamolodchikov $\mathcal{W}_3$ algebra has the same infinite-tower property
as the Virasoro algebra: each arity produces new obstructions from the
nonlinear composite $\Lambda$.
\end{proof}

\begin{remark}[Holographic complexity hierarchy]
% label removed: rem:holographic-complexity-hierarchy
The shadow depth classification produces a natural complexity hierarchy for holographic
models:
\[
\text{G (free)} \subset \text{L (current algebra)} \subset \text{C (contact)}
\subset \text{M (mixed/infinite)}.
\]
Free holographic models (such as abelian BF) have trivial modular content.
Current-algebra models (such as Kodaira--Spencer with large-$N$ gauge group) have
finite modular content determined by Lie-theoretic data.  Contact models (such as
$\beta\gamma$ holography) have finite modular content with an additional quartic
invariant.  Mixed models (such as $\mathcal{W}_3$ holography) have infinite modular
content requiring the full machinery of the shadow tower.

This hierarchy parallels the physical distinction between free, gauge, and
gravitational holographic sectors: free sectors are Gaussian, gauge sectors are
Lie-theoretic, and gravitational sectors are infinite-tower.  The shadow depth
is the algebraic invariant that quantifies this physical distinction.
\end{remark}

\subsection{The modular $R$-matrix for BF}
% label removed: subsec:modular-r-matrix-bf

\begin{definition}[Genus-zero $R$-matrix for BF]
% label removed: def:genus-zero-r-matrix-bf
The genus-zero $R$-matrix of the holomorphic BF boundary algebra is the
Maurer--Cartan kernel
\begin{equation}
% label removed: eq:r-matrix-bf
r_{\mathrm{BF}}(z) = \frac{\Omega}{z},
\end{equation}
where $\Omega = \sum_{p,q \geq 0} a_{p,q} \otimes b_{p,q}$ is the Casimir element
of the cyclic pairing.  For the abelian BF model, $\Omega$ is the identity operator
on the polynomial ring $\C[w_1, w_2]$.
\end{definition}

\begin{proposition}[Classical Yang--Baxter for BF; \ClaimStatusProvedHere]
% label removed: prop:cybe-bf
The genus-zero $R$-matrix $r_{\mathrm{BF}}(z)$ satisfies the classical Yang--Baxter
equation
\begin{equation}
% label removed: eq:cybe-bf
[r_{12}(z_{12}), r_{13}(z_{13})] + [r_{12}(z_{12}), r_{23}(z_{23})]
+ [r_{13}(z_{13}), r_{23}(z_{23})] = 0,
\end{equation}
where $z_{ij} = z_i - z_j$ and $r_{ij}$ acts on the $i$-th and $j$-th tensor
factors.
\end{proposition}

\begin{proof}
The $R$-matrix $r(z) = \Omega/z$ has the standard rational form.  The CYBE is
equivalent to the Jacobi identity for the Lie bracket dual to the cyclic pairing:
$[\Omega_{12}, \Omega_{13}] + [\Omega_{12}, \Omega_{23}] + [\Omega_{13}, \Omega_{23}] = 0$
holds in $\End(A^{\otimes 3})$ whenever $\Omega$ is the Casimir of an invariant
pairing.  For the abelian BF model, the Lie bracket is trivial, so the CYBE is
satisfied trivially.  For non-abelian BF, it follows from the Jacobi identity of the
gauge Lie algebra.
\end{proof}

\begin{theorem}[Modular $R$-matrix for abelian BF; \ClaimStatusProvedHere]
% label removed: thm:modular-r-matrix-bf
For the abelian holomorphic BF model on $S^3$, the modular $R$-matrix is trivial:
\begin{equation}
% label removed: eq:modular-r-matrix-bf
R_{\mathrm{BF}}^{\mathrm{mod}}(z; \h) = r_{\mathrm{BF}}(z) = \frac{\Omega}{z}.
\end{equation}
All genus corrections vanish: $r_g(z) = 0$ for $g \geq 1$.  This is consistent with
$\kappa = 0$ and the class~G shadow depth.
\end{theorem}

\begin{proof}
The genus-$g$ correction $r_g(z)$ to the $R$-matrix is the arity-$2$ component of
$\Theta_\cA^{(g)}$, which is $\kappa \cdot \eta \otimes \Lambda_g$ by the universal
formula.  Since $\kappa = 0$, all genus corrections vanish, and the modular $R$-matrix
equals the classical one.
\end{proof}

\begin{remark}[Non-abelian modular $R$-matrix]
% label removed: rem:non-abelian-modular-r
For non-abelian BF with gauge group $G$, the modular $R$-matrix has nontrivial genus
corrections:
\[
R^{\mathrm{mod}}_{G}(z; \h) = \frac{\Omega_G}{z} + \sum_{g \geq 1} \h^{2g}\,r_g(z),
\]
where $r_g(z) = \kappa(G) \cdot \Lambda_g \cdot \Omega_G / z$ by the universal formula.
The modular $R$-matrix is therefore a rescaled version of the classical one, with
$\h$-dependent normalization:
\[
R^{\mathrm{mod}}_G(z; \h)
= \left(1 + \sum_{g \geq 1} \h^{2g}\, \kappa(G) \Lambda_g\right)
\cdot \frac{\Omega_G}{z}.
\]
The series in parentheses is the same generating function that determines the partition
function, confirming the consistency of the modular $R$-matrix with the genus expansion.
\end{remark}

\subsection{Synthesis}
% label removed: subsec:bf-synthesis

\begin{remark}[Summary of BF computations]
% label removed: rem:bf-summary
The higher-vertex calculus for BF on $S^3$ yields the following picture:
\begin{enumerate}[label=\textup{(\arabic*)}]
\item The abelian BF model is class~G with $\kappa = 0$, hence all modular invariants
  vanish.  The model lifts trivially to all genera.
\item The non-abelian BF model is class~L, and the full modular data is determined by
  the gauge group $G$: the level $k$, rank $\dim G$, and structure constants $f^c_{ab}$.
\item The one-vertex calculus (arities $2$--$5$) produces a rigid support structure on
  the lowest modes, with unique internal test modes and finite algebraic contractions.
\item The multi-vertex wheels exhibit a matrix-trace structure: the $k$-vertex amplitude
  is a cyclic trace of a product of the cubic coefficient matrix and the propagator matrix.
\item The convergence of the one-wheel sum is controlled by the spectral radius of the
  loop transfer operator, leading to a determinant formula for the modular class.
\item The shadow depth classification of the three holographic atoms
  ($\beta\gamma$ = C, KS = L, $\mathcal{W}_3$ = M) captures the physical complexity
  hierarchy: free, gauge, gravitational.
\end{enumerate}
This is the sharpest available account of the BF model within the modular Koszul
framework.  The charged modular jet $\mathfrak{J}^{\mathrm{low}}_{\mathrm{ch}} =
(Q_{110}, Q_{011}, B_{\mathrm{mix}})$ of Definition~\ref{def:low-mode-charged-modular-jet}
is the first nonlinear invariant of the boundary algebra, and its vanishing for
abelian BF ($Q_{110} = Q_{011} = B_{\mathrm{mix}} = 0$ when $\kappa = 0$) is the
modular shadow of the trivial bulk partition function $Z_{\mathrm{BF}}(S^3) = 1$.
\end{remark}
